\documentclass[12pt]{article}

\usepackage{amsmath,amssymb,amsfonts,mathtools}
\usepackage{physics}
\usepackage{bm}
\usepackage{geometry}
\usepackage{graphicx}
\usepackage{hyperref}
\usepackage{mathrsfs}
\usepackage{tensor}
\usepackage{bbm}
\usepackage{cite}
\usepackage{setspace}
\usepackage{titlesec}

\geometry{margin=1in}
\setstretch{1.15}

\title{\textbf{Axial Holonomy and Visibility-Weighted Coherence: \\
Cosmic Birefringence in the Relativistic Scalar--Vector Plenum}}

\author{Flyxion}

\date{\today}

\begin{document}

\maketitle

\begin{abstract}

Cosmic birefringence, typically described through a Chern--Simons coupling between photons and axion-like particles, is reinterpreted here within the Relativistic Scalar--Vector Plenum (RSVP) as a visibility-weighted axial holonomy of a parity-odd sector of the plenum connection. The polarization rotation angle is formulated as the holonomy of an effective optical connection along null geodesics, while the observed washout bound is shown to constrain the phase variance of this axial sector across the finite thickness of the recombination visibility kernel.

We derive a transport formalism in which mean rotation and depolarization are governed respectively by the first and second spectral moments of the axial connection. The two-field axion mechanism of Shao--Obata--Zhang emerges as a special case of a multi-modal axial sector whose independent phase components factorize under visibility averaging, yielding the familiar $\sqrt{N}$ relaxation of washout constraints through variance additivity. A continuous synchronization model for the parity-odd phase field leads to a coherence theorem: the CMB polarization amplitude imposes a lower bound on the axial correlation length at $z \approx 1100$.

This reframing replaces particle-specific exclusion statements with a geometric coherence inequality and reveals a structural equivalence between cosmological birefringence and phase-stable photonic computation, both governed by entropy-limited holonomy across a finite integration kernel. The geometric interpretation further predicts discriminants absent in minimal axion models, including correlations between birefringence anisotropies and entropy gradients, vorticity statistics, and multipole-dependent depolarization. Cosmic birefringence is thus recast as a probe of axial synchronization in the universe’s parity-odd sector rather than solely as evidence for a propagating pseudoscalar field.

\end{abstract}

\newpage
\section{Introduction}

Recent measurements of cosmic microwave background (CMB) polarization have revived interest in the possibility of a nonzero global birefringence angle. Conventionally, such rotation is modeled through a Chern--Simons coupling of the electromagnetic field to a pseudoscalar degree of freedom, leading to a boundary-term expression for the polarization angle proportional to the displacement of an axion-like field between emission and observation. Within that framework, the observational bound on polarization amplitude constrains the oscillation rate and coupling of the pseudoscalar through a washout condition arising from the finite thickness of recombination.

While this particle-based interpretation is algebraically sufficient, it is not logically necessary. The observable quantity is a rotation of the photon’s transverse polarization frame accumulated along null geodesics and weighted by the visibility function of recombination. At a structural level, this is a holonomy problem: polarization is parallel-transported by an effective optical connection, and birefringence measures the parity-odd component of that connection integrated along the path. The washout bound does not constrain a field species per se; it constrains the phase coherence of whatever geometric sector sources that parity-odd projection across the emission kernel.

In this work we develop that geometric interpretation within the Relativistic Scalar--Vector Plenum (RSVP) framework. RSVP posits that cosmological dynamics are governed not solely by a metric tensor but by a coupled scalar--vector--entropy system $(\Phi, v^\mu, S)$ defined on spacetime. The scalar field $\Phi$ encodes plenum density or potential structure, the vector field $v^\mu$ represents organized flow, and the entropy field $S$ encodes thermodynamic ordering. Gravitational and optical phenomena emerge from the combined geometry of these fields through effective connections constructed from their gradients and contractions. In particular, parity-odd structures arise naturally from pseudoscalar combinations such as $\varepsilon_{\mu\nu\rho\sigma} k^\mu u^\nu \nabla^\rho S v^\sigma$, where $k^\mu$ is the photon wavevector and $u^\mu$ a comoving four-velocity.

We show that cosmic birefringence can be formulated as the axial holonomy of such a parity-odd connection. The observable rotation angle is not the raw line integral of the axial projection but its convolution with the recombination visibility function. The empirical washout constraint therefore becomes a bound on the phase variance of the parity-odd sector across the recombination thickness. This leads to a coherence theorem: the axial correlation length of the plenum at $z \approx 1100$ must exceed the visibility width in order to produce rotation without depolarization.

To make this coherence constraint explicit, we introduce a synchronization model for the parity-odd phase field. Beginning with a discrete five-dimensional Ising representation and passing to a continuous axial phase description, we demonstrate that the familiar two-field relaxation of washout constraints follows generically from variance additivity under independent phase modes. The mechanism does not depend on the existence of multiple particle species; it is a property of multi-modal axial coherence.

This reinterpretation reveals a deeper structural parallel between cosmological birefringence and engineered phase-based computation. In both cases, coherent holonomy is accumulated across a finite integration kernel, and the preservation of amplitude is limited by phase diffusion. The cosmological washout bound is therefore mathematically equivalent to an entropy--coherence inequality governing photonic and thermodynamic computing systems.

Finally, we derive observational discriminants capable of distinguishing geometric axial coherence from minimally coupled pseudoscalar transport. While the two interpretations coincide at the level of isotropic mean rotation, they differ in predicted cross-correlations, multipole-dependent depolarization, and higher-order phase statistics. These differences render the geometric interpretation empirically testable.

The paper proceeds by first deriving the axial transport law in the RSVP optical connection, then establishing the visibility-weighted holonomy formalism and coherence theorem. We subsequently develop the synchronization model and entropy--coherence inequality, and conclude by identifying observational signatures that differentiate geometric parity-odd coherence from free pseudoscalar propagation.

\newpage
\section{Geometric Structure of the Relativistic Scalar--Vector Plenum}

The Relativistic Scalar--Vector Plenum (RSVP) is a geometric field theory in which spacetime is not treated as an empty background manifold endowed solely with a metric structure, but as a dynamically ordered medium characterized by three interlocking fields: a scalar density field $\Phi$, a vector flow field $v^\mu$, and an entropy density field $S$. The fundamental ontology is that of a structured plenum whose macroscopic geometric properties arise from the coupled dynamics of these fields rather than from vacuum curvature alone.

The scalar field $\Phi$ encodes the compressional state of the plenum and controls effective volumetric response. The vector field $v^\mu$ defines directed transport and encodes large-scale ordering and flow coherence. The entropy field $S$ parameterizes the degree of local disorder and governs irreversible relaxation dynamics. The metric $g_{\mu\nu}$ is not eliminated, but it is interpreted as an emergent large-scale descriptor of averaged plenum structure rather than a primitive background variable. In particular, deviations from Levi--Civita transport arise from gradients and vorticity of $(\Phi, v^\mu, S)$.

The RSVP connection is therefore not restricted to the torsion-free Levi--Civita connection $\Gamma^\mu_{\nu\rho}$ derived from $g_{\mu\nu}$, but admits corrections generated by entropy gradients and vector-field structure. We write the full plenum connection as
\begin{equation}
\widetilde{\Gamma}^\mu_{\ \nu\rho}
=
\Gamma^\mu_{\ \nu\rho}
+
\Delta^\mu_{\ \nu\rho}(\Phi, v, S),
\end{equation}
where $\Delta^\mu_{\ \nu\rho}$ represents non-Riemannian structure induced by plenum ordering.

The correction term $\Delta^\mu_{\ \nu\rho}$ may be decomposed irreducibly into parity-even and parity-odd components under spatial inversion. The parity-even sector modifies effective refractive and dispersive properties of null transport without inducing handedness-sensitive effects. The parity-odd sector, however, acts differently on left- and right-handed polarization states. This sector is the geometric locus of cosmic birefringence in the RSVP framework.

Let $k^\mu = \frac{dx^\mu}{d\lambda}$ denote the null tangent to a photon worldline $\gamma$. Let $P^\mu$ be a transverse polarization vector satisfying $P^\mu k_\mu = 0$ and normalized within the screen space orthogonal to $k^\mu$ and the observer four-velocity $u^\mu$. In standard geometric optics, polarization is parallel transported via
\begin{equation}
k^\nu \nabla_\nu P^\mu = 0.
\end{equation}
In RSVP, polarization transport is governed instead by
\begin{equation}
k^\nu \widetilde{\nabla}_\nu P^\mu
=
k^\nu \nabla_\nu P^\mu
+
k^\nu \Delta^\mu_{\ \nu\rho} P^\rho
=
0.
\end{equation}

Projecting onto the two-dimensional screen bundle orthogonal to both $k^\mu$ and $u^\mu$, the transport equation reduces to an $SO(2)$ rotation in the polarization plane. The generator of such rotations may be written as $\epsilon^{ab}$ in the orthonormal polarization basis $\{e_1^a, e_2^a\}$. The parity-odd part of $\Delta^\mu_{\ \nu\rho}$ induces an effective scalar angular velocity $\kappa_A$ along the null path:
\begin{equation}
\frac{d\theta}{d\lambda} = \kappa_A(\lambda),
\end{equation}
where $\theta$ is the polarization angle and
\begin{equation}
\kappa_A(\lambda)
=
\Pi^\mu_{\ a}
\Pi^\nu_{\ b}
\epsilon^{ab}
\, k^\rho
\Delta_{\mu\nu\rho}^{(\text{odd})}.
\end{equation}
Here $\Pi^\mu_{\ a}$ projects onto the screen space and $\Delta^{(\text{odd})}$ denotes the parity-odd component of the plenum connection.

The observable polarization rotation accumulated along the null path $\gamma$ from emission at conformal time $\eta_{\rm LSS}$ to observation at $\eta_0$ is therefore the axial holonomy
\begin{equation}
H_A[\gamma]
=
\int_{\eta_{\rm LSS}}^{\eta_0}
\kappa_A(\eta, x(\eta)) \, d\eta.
\end{equation}

This expression is structurally identical to the Chern--Simons axion result
\begin{equation}
\beta
=
\frac{1}{2} g_{\phi\gamma}
\int_{\eta_{\rm LSS}}^{\eta_0}
\partial_\eta \phi \, d\eta,
\end{equation}
provided the identification
\begin{equation}
\kappa_A
\;\longleftrightarrow\;
\frac{1}{2} g_{\phi\gamma} \partial_\eta \phi
\end{equation}
is made.

In the axion model, $\phi$ is a propagating pseudoscalar particle species. In RSVP, $\kappa_A$ is instead a projection of geometric structure. The pseudoscalar degree of freedom becomes a coordinate chart on the parity-odd sector of $\widetilde{\Gamma}$ rather than an ontological primitive.

The significance of this reinterpretation is immediate. The birefringence angle $\beta$ is not a field displacement but a holonomy of the plenum's axial connection. The question becomes not what particle causes $\beta$, but what geometric coherence properties of the parity-odd sector permit a non-vanishing holonomy to survive the recombination visibility averaging.

Before turning to visibility weighting and washout, we clarify the internal RSVP origin of $\Delta^{(\text{odd})}$. A minimal parity-odd contribution consistent with covariance may be constructed from the totally antisymmetric Levi--Civita tensor $\varepsilon_{\mu\nu\rho\sigma}$ and a pseudoscalar density $\Xi$ built from plenum fields. For example,
\begin{equation}
\Delta^{(\text{odd})}_{\mu\nu\rho}
=
\alpha
\, \varepsilon_{\mu\nu\rho\sigma}
\, \nabla^\sigma \Xi(\Phi, v, S),
\end{equation}
where $\alpha$ is a coupling scale and $\Xi$ is any pseudoscalar functional of the RSVP fields, such as
\begin{equation}
\Xi
=
\nabla_\mu S \, \omega^\mu,
\qquad
\omega^\mu
=
\varepsilon^{\mu\nu\rho\sigma}
u_\nu \nabla_\rho v_\sigma,
\end{equation}
which couples entropy gradients to vorticity and is manifestly parity-odd.

In this construction, cosmic birefringence arises from large-scale coherence in entropy–vorticity coupling. The axion model is recovered as a special case where $\Xi$ is replaced by a fundamental scalar field $\phi$ with its own dynamics. RSVP instead treats $\Xi$ as emergent from plenum ordering.

We now proceed to incorporate the finite thickness of last scattering and derive the visibility-weighted axial holonomy, which produces the washout effect as a spectral coherence constraint on $\kappa_A$.

\newpage
\section{Visibility-Weighted Holonomy and Coherence Constraints}

The observable cosmic microwave background polarization is not generated at a sharply defined hypersurface but across a finite conformal time interval centered at recombination. Let $\eta$ denote conformal time and $W(\eta)$ the visibility function, normalized so that
\begin{equation}
\int_{-\infty}^{\infty} W(\eta)\, d\eta = 1,
\end{equation}
with $W(\eta)$ sharply peaked near $\eta_\star$ corresponding to redshift $z \approx 1100$, and characteristic width $\Delta\eta_{\mathrm{rec}}$.

Let $P_\pm = Q \pm iU$ denote the complex polarization amplitudes transforming under a rotation of the polarization angle $\theta$ as
\begin{equation}
P_\pm \mapsto e^{\pm 2 i \theta} P_\pm.
\end{equation}
In the absence of parity-odd structure, the emitted polarization $P_\pm^{(0)}(\eta,\hat n)$ is transported trivially along the null ray to the observer. In RSVP, the axial holonomy modifies this transport so that the observed polarization is
\begin{equation}
P_\pm^{\mathrm{obs}}(\hat n)
=
\int d\eta \,
W(\eta)
\, e^{\pm 2 i \beta(\eta,\hat n)}
\, P_\pm^{(0)}(\eta,\hat n),
\end{equation}
where
\begin{equation}
\beta(\eta,\hat n)
=
\int_{\eta}^{\eta_0}
\kappa_A(\eta',x(\eta'))\, d\eta'
\end{equation}
is the axial holonomy accumulated from emission time $\eta$ to the observer at $\eta_0$.

Assuming that the intrinsic polarization varies slowly across the narrow support of $W(\eta)$, we may factor it out of the integral to leading order, yielding
\begin{equation}
P_\pm^{\mathrm{obs}}(\hat n)
=
\mathcal{W}_\pm(\hat n)
\, P_\pm^{(0)}(\eta_\star,\hat n),
\end{equation}
with the visibility-weighted phase factor
\begin{equation}
\mathcal{W}_\pm(\hat n)
=
\int d\eta \,
W(\eta)
\, e^{\pm 2 i \beta(\eta,\hat n)}.
\end{equation}

The observable polarization amplitude is therefore suppressed relative to its emission value by
\begin{equation}
\left|\mathcal{W}_\pm\right|.
\end{equation}
The washout effect is precisely the condition $\left|\mathcal{W}_\pm\right| < 1$ arising from phase decoherence across the emission kernel.

To expose the coherence structure, decompose the axial connection projection into spectral modes along the null path:
\begin{equation}
\kappa_A(\eta)
=
\kappa_0(\eta)
+
\sum_{k}
a_k \cos(\omega_k \eta + \varphi_k),
\end{equation}
where $\kappa_0$ is the slowly varying component and the oscillatory modes represent finite-frequency structure in the parity-odd sector of the plenum connection.

The holonomy from $\eta$ to $\eta_0$ is then
\begin{equation}
\beta(\eta)
=
\bar{\beta}
+
\sum_k
A_k
\sin(\omega_k \eta + \varphi_k),
\end{equation}
where
\begin{equation}
A_k = \frac{a_k}{\omega_k}
\end{equation}
and
\begin{equation}
\bar{\beta}
=
\int_{\eta_\star}^{\eta_0}
\kappa_0(\eta') d\eta'
\end{equation}
denotes the coherent component accumulated outside the rapidly oscillating regime.

Substituting into $\mathcal{W}_\pm$ gives
\begin{equation}
\mathcal{W}_\pm
=
e^{\pm 2 i \bar{\beta}}
\int d\eta \,
W(\eta)
\,
\exp
\left[
\pm 2 i
\sum_k
A_k
\sin(\omega_k \eta + \varphi_k)
\right].
\end{equation}

When $\omega_k \Delta\eta_{\mathrm{rec}} \gg 1$, the oscillations are rapid relative to the kernel width and the integral reduces to a Bessel suppression factor:
\begin{equation}
\int d\eta \,
W(\eta)
\, e^{i A_k \sin(\omega_k \eta)}
\approx
J_0(A_k).
\end{equation}

If the oscillatory modes are spectrally separated such that cross terms average to zero, the suppression factor factorizes:
\begin{equation}
\mathcal{W}_\pm
\approx
e^{\pm 2 i \bar{\beta}}
\prod_k
J_0(2 A_k).
\end{equation}

The mean rotation angle is therefore controlled by $\bar{\beta}$, while the polarization amplitude is reduced by the product of Bessel factors. Expanding for small amplitudes,
\begin{equation}
J_0(2 A_k)
=
1 - 2 A_k^2 + \mathcal{O}(A_k^4),
\end{equation}
so that
\begin{equation}
\left|\mathcal{W}_\pm\right|
\approx
1 - 2 \sum_k A_k^2 + \mathcal{O}(A_k^4).
\end{equation}

We therefore obtain the fundamental coherence inequality:
\begin{equation}
\sum_k A_k^2
\lesssim
\epsilon_{\mathrm{wash}},
\end{equation}
where $\epsilon_{\mathrm{wash}}$ is the observationally determined upper bound on depolarization.

This relation shows that the washout constraint is a bound on the spectral variance of the axial holonomy integrand at frequencies exceeding $\Delta\eta_{\mathrm{rec}}^{-1}$. It does not constrain the mean coherent component $\bar{\beta}$ directly.

The separation between coherent and incoherent modes is therefore defined by the dimensionless parameter
\begin{equation}
\omega_k \Delta\eta_{\mathrm{rec}}.
\end{equation}
Modes satisfying
\begin{equation}
\omega_k \Delta\eta_{\mathrm{rec}} \ll 1
\end{equation}
remain coherent across the visibility kernel and contribute to $\bar{\beta}$ without washout suppression. Modes satisfying
\begin{equation}
\omega_k \Delta\eta_{\mathrm{rec}} \gg 1
\end{equation}
contribute negligibly to the mean rotation but contribute quadratically to depolarization.

This establishes the central reinterpretation: the empirical washout bound is a lower bound on the coherence length of the parity-odd sector of the plenum connection at recombination. It is a spectral constraint on $\kappa_A$, not a particle-exclusion constraint.

We now demonstrate that the two-field axion mechanism is a special case of this general multi-modal structure.

\newpage
\section{Multi-Modal Axial Connection and Factorized Dephasing}

We now make explicit the correspondence between the two-field axion scenario and the general multi-modal axial connection derived above. In the conventional Chern--Simons axion model, the axial rotation rate is proportional to the time derivative of a pseudoscalar field,
\begin{equation}
\kappa_A(\eta)
=
\frac{1}{2} g_{\phi\gamma} \partial_\eta \phi(\eta).
\end{equation}
If two independent axion-like fields $\phi_1$ and $\phi_2$ are present with identical photon couplings, one has
\begin{equation}
\kappa_A(\eta)
=
\frac{1}{2} g_{\phi\gamma}
\left(
\partial_\eta \phi_1(\eta)
+
\partial_\eta \phi_2(\eta)
\right).
\end{equation}

Assuming oscillatory solutions during matter domination of the form
\begin{equation}
\phi_i(\eta)
=
\Phi_i \cos(m_i \eta + \alpha_i),
\end{equation}
the axial projection becomes
\begin{equation}
\kappa_A(\eta)
=
-
\frac{1}{2} g_{\phi\gamma}
\sum_{i=1}^2
\Phi_i m_i
\sin(m_i \eta + \alpha_i).
\end{equation}

Comparing with the general RSVP expansion
\begin{equation}
\kappa_A(\eta)
=
\sum_k a_k \cos(\omega_k \eta + \varphi_k),
\end{equation}
we identify
\begin{equation}
a_i
=
\frac{1}{2} g_{\phi\gamma} \Phi_i m_i,
\qquad
\omega_i = m_i.
\end{equation}

The axial holonomy contribution from each mode is then
\begin{equation}
A_i
=
\frac{a_i}{\omega_i}
=
\frac{1}{2} g_{\phi\gamma} \Phi_i.
\end{equation}

Crucially, the amplitude entering the washout suppression depends only on the field amplitude $\Phi_i$, not on $m_i$ directly. The mass controls instead the coherence classification through the dimensionless product $m_i \Delta\eta_{\mathrm{rec}}$.

Inserting two modes into the visibility-weighted factor,
\begin{equation}
\mathcal{W}_\pm
=
e^{\pm 2 i \bar{\beta}}
\int d\eta \,
W(\eta)
\,
\exp
\left[
\pm 2 i
\left(
A_1 \sin(m_1 \eta + \alpha_1)
+
A_2 \sin(m_2 \eta + \alpha_2)
\right)
\right],
\end{equation}
and assuming $|m_1 - m_2| \gg \Delta\eta_{\mathrm{rec}}^{-1}$ so that cross terms average out, one obtains
\begin{equation}
\mathcal{W}_\pm
\approx
e^{\pm 2 i \bar{\beta}}
J_0(2 A_1)
J_0(2 A_2).
\end{equation}

This is precisely the factorization structure exploited in the two-field analysis of cosmic birefringence. The depolarization factor becomes a product of Bessel functions rather than a single Bessel function evaluated at the sum amplitude.

To understand why this relaxes the washout constraint, consider small amplitudes $A_i \ll 1$. For a single effective mode with total amplitude $A = A_1 + A_2$, the suppression would scale as
\begin{equation}
J_0(2A)
\approx
1 - 2(A_1 + A_2)^2.
\end{equation}
By contrast, the factorized two-mode suppression gives
\begin{equation}
J_0(2A_1)J_0(2A_2)
\approx
1 - 2(A_1^2 + A_2^2).
\end{equation}

The difference between $(A_1 + A_2)^2$ and $A_1^2 + A_2^2$ is the cross term $2A_1A_2$. In the single-mode interpretation, coherent addition of amplitudes increases depolarization quadratically in the total amplitude. In the multi-modal case with independent oscillatory phases across the visibility kernel, this cross term does not contribute to washout. The suppression is therefore weaker for fixed total coherent rotation.

The mean rotation angle observed at the detector, however, is controlled by the boundary contribution
\begin{equation}
\beta_{\mathrm{mean}}
=
\frac{1}{2} g_{\phi\gamma}
\left(
\phi_1(\eta_0)
+
\phi_2(\eta_0)
\right),
\end{equation}
provided the emission contribution averages away. The coherent mean rotation thus scales linearly with $\Phi_1 + \Phi_2$, while the washout constraint scales with $\Phi_1^2 + \Phi_2^2$.

At equal amplitudes $\Phi_1 = \Phi_2$, the ratio of allowed coherent rotation to washout variance increases by a factor of $\sqrt{2}$ relative to the single-field case. This is not a special property of axions but a generic consequence of variance additivity under factorized oscillatory modes.

Within RSVP, this result is interpreted as follows. The parity-odd plenum connection may possess multiple oscillatory modes with incommensurate coherence lengths during recombination. The visibility-weighted dephasing factorizes across such modes. The observational constraint on depolarization therefore bounds the sum of squared mode amplitudes rather than the square of their coherent sum.

The two-field axion mechanism is thus recognized as a special instance of a general geometric statement: multi-modal axial holonomy can raise the achievable coherent polarization rotation at fixed dephasing budget.

It is important to emphasize that this mechanism depends only on spectral separation across the visibility kernel. If $|m_1 - m_2| \lesssim \Delta\eta_{\mathrm{rec}}^{-1}$, the cross terms do not average out and the factorization fails. In RSVP language, this corresponds to two parity-odd modes whose coherence scales overlap at recombination, effectively behaving as a single composite mode.

We now turn to the question of local-field dominance and the structural origin of the boundary-term form of the observable rotation.

\newpage
\section{Boundary Dominance and Foliation Structure of Axial Holonomy}

The standard axion derivation of cosmic birefringence yields a boundary-term expression,
\begin{equation}
\beta
=
\frac{1}{2} g_{\phi\gamma}
\left(
\phi(\eta_0) - \phi(\eta_{\mathrm{LSS}})
\right),
\end{equation}
suggesting that the observable rotation depends only on field values at emission and reception. In the regime relevant for ultralight oscillatory modes, the emission contribution averages to zero over the finite recombination thickness, leaving the observer-side value as the dominant term.

Within RSVP, this boundary dominance is not an accident of the pseudoscalar equation of motion but a structural consequence of visibility-weighted holonomy.

Recall that the axial holonomy accumulated from emission time $\eta$ to the observer is
\begin{equation}
\beta(\eta)
=
\int_{\eta}^{\eta_0}
\kappa_A(\eta')\, d\eta'.
\end{equation}
The observable rotation is the visibility-weighted average
\begin{equation}
\beta_{\mathrm{obs}}
=
\int d\eta\, W(\eta)\, \beta(\eta).
\end{equation}
Substituting,
\begin{equation}
\beta_{\mathrm{obs}}
=
\int d\eta\, W(\eta)
\int_{\eta}^{\eta_0}
\kappa_A(\eta')\, d\eta'.
\end{equation}

Interchanging integration order,
\begin{equation}
\beta_{\mathrm{obs}}
=
\int_{\eta_{\mathrm{LSS}}}^{\eta_0}
\kappa_A(\eta')
\left(
\int_{\eta_{\mathrm{LSS}}}^{\eta'} W(\eta)\, d\eta
\right)
d\eta'.
\end{equation}

Define the cumulative visibility function
\begin{equation}
\mathcal{F}(\eta')
=
\int_{\eta_{\mathrm{LSS}}}^{\eta'} W(\eta)\, d\eta.
\end{equation}
Then
\begin{equation}
\beta_{\mathrm{obs}}
=
\int_{\eta_{\mathrm{LSS}}}^{\eta_0}
\kappa_A(\eta')\, \mathcal{F}(\eta')\, d\eta'.
\end{equation}

The function $\mathcal{F}(\eta')$ rises sharply from $0$ to $1$ across the narrow width $\Delta\eta_{\mathrm{rec}}$ and remains unity thereafter. Consequently,
\begin{equation}
\beta_{\mathrm{obs}}
=
\int_{\eta_\star}^{\eta_0}
\kappa_A(\eta')\, d\eta'
+
\int_{\eta_{\mathrm{LSS}}}^{\eta_\star}
\kappa_A(\eta')\, \mathcal{F}(\eta')\, d\eta'.
\end{equation}

The second term involves integration over the narrow recombination window weighted by $\mathcal{F}(\eta') \in (0,1)$. For oscillatory modes satisfying $\omega_k \Delta\eta_{\mathrm{rec}} \gg 1$, this contribution averages to zero. The first term extends from $\eta_\star$ to $\eta_0$ without kernel suppression.

We therefore obtain
\begin{equation}
\beta_{\mathrm{obs}}
\approx
\int_{\eta_\star}^{\eta_0}
\kappa_A(\eta')\, d\eta',
\end{equation}
provided the high-frequency emission-window contribution vanishes.

This expression makes clear that the observable rotation is the difference between the axial state of the plenum evaluated at two foliations: the recombination foliation and the observer foliation. If $\kappa_A$ itself oscillates many times between $\eta_\star$ and $\eta_0$, then the value at $\eta_\star$ behaves effectively as a random phase across the ensemble of modes contributing to the local plenum state. Its contribution vanishes upon ensemble averaging, leaving
\begin{equation}
\beta_{\mathrm{obs}}
\approx
\int^{\eta_0} \kappa_A(\eta')\, d\eta'
-
\int^{\eta_\star} \kappa_A(\eta')\, d\eta'
\;\longrightarrow\;
\int^{\eta_0} \kappa_A(\eta')\, d\eta',
\end{equation}
in the statistical sense.

This is the geometric origin of local-field dominance. The observer foliation corresponds to a single evaluation of the plenum state, whereas the emission foliation samples an ensemble of oscillatory phases across the recombination thickness. The latter contributes variance but not coherent mean.

Formally, if $\kappa_A$ is decomposed into modes with random phases at $\eta_\star$, then
\begin{equation}
\left\langle
\int_{\eta_{\mathrm{LSS}}}^{\eta_\star}
\kappa_A(\eta')\, \mathcal{F}(\eta')\, d\eta'
\right\rangle
=
0,
\end{equation}
while
\begin{equation}
\int_{\eta_\star}^{\eta_0}
\kappa_A(\eta')\, d\eta'
=
\sum_k
A_k
\sin(\omega_k \eta_0 + \varphi_k),
\end{equation}
which depends only on the local plenum phase.

The observed isotropic birefringence angle therefore measures the axial state of the plenum at the observer's foliation, subject to coherence constraints inherited from recombination.

This clarifies the reinterpretation of the axion boundary term. The apparent field displacement $\phi(\eta_0) - \phi(\eta_{\mathrm{LSS}})$ is, in RSVP language, the difference in axial plenum ordering between two spacetime slices. The washout effect ensures that only those spectral components of $\kappa_A$ that are coherent across the recombination thickness contribute to this difference as mean rotation rather than as depolarizing noise.

We may now reformulate the preceding transport analysis as a general coherence statement for the parity-odd sector of the plenum. The visibility-weighted holonomy formalism established above shows that the observable birefringence angle depends on the axial connection only through its phase variance across the finite recombination kernel. The empirical washout bound therefore constrains not the magnitude of $\kappa_A$ alone, but the correlation structure of the underlying parity-odd degrees of freedom along the emission foliation.

More precisely, the previous section implies that any axial sector whose correlation length $\xi_{\mathrm{rec}}$ satisfies
\begin{equation}
\xi_{\mathrm{rec}} \gg \Delta\eta_{\mathrm{rec}}
\end{equation}
contributes coherently to mean rotation, whereas sectors with
\begin{equation}
\xi_{\mathrm{rec}} \ll \Delta\eta_{\mathrm{rec}}
\end{equation}
contribute predominantly to depolarization. The observational bound on polarization amplitude is therefore equivalent to a lower bound on the axial correlation length of the parity-odd plenum sector at recombination.

This statement is independent of microphysical realization. It is a geometric consequence of visibility-weighted holonomy. The remaining question is constructive: what class of mesoscopic dynamics can generate an axial sector with tunable correlation length and controlled phase variance?

To answer this, we introduce a synchronization model that provides an explicit realization of the required coherence structure.

\newpage
\section{A 5D Ising Synchronization Model for Axial Coherence and Its Computing Implications}

The reinterpretation of washout as a coherence constraint suggests that one should not begin with a particle spectrum but with a mesoscopic dynamics for the parity-odd plenum sector near recombination. The essential empirical fact is kernel-limited: polarization transport is visibility-weighted, hence any axial degree of freedom whose coherence length along the emission foliation is shorter than the recombination thickness contributes predominantly to dephasing rather than to mean rotation. This is a statement about synchronization across a finite interval, and it admits a natural lattice realization in which the parity-odd sector is modeled as an order parameter with a tunable correlation length.

To that end we introduce a five-dimensional Ising-type synchronization model. The use of five dimensions is not a claim about embedding spacetime in a literal higher-dimensional manifold, but a device for representing an internal axial mode as an additional lattice direction, thereby enabling one to treat coherence in that sector as a bona fide correlation-length problem rather than a purely kinematic average. Concretely, let $\Lambda \subset \mathbb{Z}^5$ be a hypercubic lattice with sites labeled by
\begin{equation}
x = (n_0,n_1,n_2,n_3,n_4),
\end{equation}
where the first four coordinates represent a coarse-grained comoving spacetime lattice in conformal gauge and $n_4$ indexes an internal axial layer. At each site we place an Ising spin
\begin{equation}
\sigma_x \in \{-1,+1\},
\end{equation}
interpreted as the local sign of an axial response variable, i.e.\ the sign of the parity-odd projection of the RSVP optical connection, or more precisely the sign of a coarse-grained surrogate for $\kappa_A$ after normalization by an amplitude scale.

We define the anisotropic Hamiltonian
\begin{equation}
\mathcal{H}[\sigma]
=
-\sum_{x\in\Lambda}
\left(
\sum_{\mu=0}^{3} J_\mu \, \sigma_x \sigma_{x+\hat\mu}
+
J_4 \, \sigma_x \sigma_{x+\hat 4}
\right)
-\sum_{x\in\Lambda} h_x \sigma_x,
\end{equation}
where $J_\mu$ controls coupling along the four spacetime directions, $J_4$ controls coupling across axial layers, and $h_x$ is an external field encoding slow RSVP biases, which may be taken to depend on $(\Phi,v,S)$ through a parity-odd pseudoscalar functional $\Xi$ as $h_x \propto \Xi(x)$. The fundamental synchronization parameter in this model is the axial correlation length $\xi_4$ induced by $J_4$ at inverse temperature $\beta_T$, which sets the scale over which the parity-odd sector remains ordered across the internal dimension.

The observable relevant for cosmic birefringence is not the microscopic spin itself but a coarse-grained axial connection variable obtained by a block-spin transform and projection onto a null ray. Let $\gamma$ be a discrete null path through the four-dimensional sublattice (holding $n_4$ fixed or averaged, depending on the observational reduction), and define a coarse axial integrand
\begin{equation}
\kappa_A(\eta)
\;\equiv\;
\kappa_0 \, m(\eta),
\qquad
m(\eta)
=
\mathbb{E}\!\left[\sigma_{x(\eta)}\right],
\end{equation}
where $\kappa_0$ sets dimensions and $m(\eta)$ is the local magnetization along the path in the appropriate ensemble. The axial holonomy from emission time $\eta$ to observation is then discretized as
\begin{equation}
\beta(\eta)
=
\sum_{\eta' \ge \eta}
\kappa_A(\eta') \Delta\eta',
\end{equation}
and the visibility-weighted observable is
\begin{equation}
\beta_{\mathrm{obs}}
=
\sum_{\eta}
W(\eta)\,\beta(\eta)\,\Delta\eta.
\end{equation}

The washout effect depends not on $\mathbb{E}[\beta]$ alone but on the characteristic function of $\beta(\eta)$ across the emission kernel. In the small-angle regime one may write the polarization suppression in terms of the variance of the axial phase across the kernel. Denoting the random variable
\begin{equation}
\Theta
=
2\beta(\eta)
\end{equation}
(with the factor of $2$ corresponding to the spin-2 nature of polarization), the leading depolarization factor is
\begin{equation}
\left|\mathcal{W}_\pm\right|
=
\left|\mathbb{E}_W\!\left[e^{\pm i\Theta}\right]\right|
\approx
\exp\!\left(-\frac{1}{2}\mathrm{Var}_W(\Theta)\right),
\end{equation}
where $\mathbb{E}_W[\cdot]$ denotes visibility-weighted averaging over emission times. Thus the observational washout bound is equivalently an upper bound on $\mathrm{Var}_W(\Theta)$.

In the 5D Ising model, the variance of $\Theta$ is controlled by the two-point correlation function of the axial order parameter across the recombination thickness. Writing $m(\eta)$ as a block observable derived from spins, one obtains generically
\begin{equation}
\mathrm{Var}_W(\Theta)
=
4\kappa_0^2
\sum_{\eta,\eta'}
W(\eta)W(\eta')
\,\mathrm{Cov}\!\left(
\sum_{\eta''\ge\eta} m(\eta'')\Delta\eta'',
\sum_{\eta'''\ge\eta'} m(\eta''')\Delta\eta'''
\right).
\end{equation}
Under stationarity across the narrow recombination window and in the regime where the dominant contributions arise from within that window, the covariance reduces to an integral transform of the connected correlation function
\begin{equation}
C(\Delta\eta)
=
\mathbb{E}\!\left[m(\eta)m(\eta+\Delta\eta)\right]-\mathbb{E}[m(\eta)]^2.
\end{equation}
If $C(\Delta\eta)$ decays exponentially with correlation length $\xi_{\mathrm{rec}}$ along the emission foliation,
\begin{equation}
C(\Delta\eta)\sim e^{-|\Delta\eta|/\xi_{\mathrm{rec}}},
\end{equation}
then the washout bound becomes, in scaling form,
\begin{equation}
\mathrm{Var}_W(\Theta)
\;\sim\;
4\kappa_0^2
\int d\eta\,d\eta'\,
W(\eta)W(\eta')\,\mathcal{K}(\eta,\eta')\,e^{-|\eta-\eta'|/\xi_{\mathrm{rec}}},
\end{equation}
for an explicit kernel $\mathcal{K}$ determined by the holonomy sums. The key qualitative regime distinction is immediate: if $\xi_{\mathrm{rec}}\gg \Delta\eta_{\mathrm{rec}}$ the exponential is near unity over the support of $W$ and the variance is suppressed relative to the mean squared, producing rotation without depolarization; if $\xi_{\mathrm{rec}}\ll \Delta\eta_{\mathrm{rec}}$ the correlations collapse and the variance becomes effectively extensive in the number of incoherent sub-intervals, producing depolarization with negligible coherent mean.

The role of the fifth dimension is now transparent. The axial correlation length $\xi_{\mathrm{rec}}$ is not postulated; it is generated by the synchronization coupling $J_4$ (and by the proximity to the critical surface in $(J_\mu,J_4,\beta_T)$-space). In particular, a near-critical parity-odd sector yields large $\xi_4$ and therefore large effective coherence along the emission foliation after projection, while a far-from-critical sector yields small $\xi_4$ and hence strong washout. The ``two-field'' or multi-modal relief mechanism becomes, in this language, the statement that $m(\eta)$ may be a sum of weakly correlated components derived from distinct axial layers or distinct coarse-grained modes. If one writes
\begin{equation}
m(\eta)=\sum_{a=1}^{N} m_a(\eta),
\end{equation}
with approximately vanishing cross-covariances on the visibility scale, then
\begin{equation}
\mathrm{Var}_W(\Theta)\propto \sum_{a=1}^{N} \mathrm{Var}_W(\Theta_a),
\end{equation}
whereas the coherent mean rotation scales as $\sum_a \mathbb{E}[\Theta_a]$. This is precisely the linear-versus-quadratic bookkeeping previously identified at the level of Bessel-factorization; the Ising synchronization model provides a microscopic mechanism for realizing the required covariance structure.

This formulation admits an immediate bridge to photonic and thermodynamic computing. The axial sector here is a synchronization resource: it is a low-entropy ordered degree of freedom whose coherence must be maintained over a kernel thickness in order to effect a controlled rotation rather than an uncontrolled depolarization. In photonic computing the analogous distinction is between coherent phase manipulation (which performs computation through controllable interference and holonomy-like transformations in mode space) and incoherent phase noise (which converts intended unitary-like transformations into amplitude loss and heating). The washout bound is therefore isomorphic to an engineering constraint: any architecture that computes by polarization holonomy must keep the axial phase variance below a threshold determined by the ``visibility kernel'' of the device, i.e.\ the temporal or spatial extent over which the output integrates the phase.

Thermodynamic computing enters because synchronization in an Ising-like medium is not free: maintaining large correlation length near criticality is paid for either by slow relaxation (critical slowing down), by energy injection, or by engineered dissipation pathways. In RSVP language this is exactly the entropy budget of geometric order. In the Ising model it is encoded by the free energy landscape: the cost of sustaining an ordered axial sector is measured by the distance from equilibrium under the imposed fields $h_x$ and by the rate at which entropy is exported to stabilize the ordered manifold. Thus the same inequality that enforces ``no washout'' in cosmology becomes, in computing, an inequality relating achievable coherent transformation depth to entropy production per operation.

In this way, a 5D Ising synchronization model provides a unifying mesoscopic picture. Cosmic birefringence becomes the large-scale holonomy produced by an ordered parity-odd sector whose coherence length exceeds the emission kernel, while washout becomes the depolarization induced by sub-kernel dephasing. Photonic computing becomes the controlled utilization of the same holonomy principle in engineered media, and thermodynamic computing becomes the accounting of the entropy required to keep the relevant sector synchronized rather than noisy. The cosmological constraint is then interpretable as a bound on allowable axial phase noise at recombination, and the engineering analogue is a bound on allowable phase noise over the integration window of the photonic device.

\newpage
\section{Continuous Axial Phase Synchronization and the Coherence Theorem}

The discrete Ising surrogate introduced above captures the combinatorial structure of synchronization but suppresses an essential geometric feature of polarization transport: the axial sector is fundamentally phase-valued rather than sign-valued. A more faithful representation of the parity-odd plenum connection therefore employs a continuous axial phase field
\begin{equation}
\theta(x) \in S^1,
\end{equation}
defined on a coarse-grained five-dimensional lattice or, in the continuum limit, on a fiber bundle over spacetime whose fiber is the axial phase manifold.

We consider an anisotropic XY-type synchronization Hamiltonian
\begin{equation}
\mathcal{H}[\theta]
=
- \sum_{x}
\left(
\sum_{\mu=0}^{3}
J_\mu
\cos(\theta_x - \theta_{x+\hat\mu})
+
J_4
\cos(\theta_x - \theta_{x+\hat 4})
\right)
- \sum_x h_x \cos(\theta_x),
\end{equation}
where the fifth-direction coupling $J_4$ controls coherence in the internal axial layer and $h_x$ encodes slow RSVP biases generated by a pseudoscalar functional
\begin{equation}
h_x
=
\lambda \, \Xi(\Phi,v,S)_x,
\qquad
\Xi
=
\nabla_\mu S \, \omega^\mu,
\end{equation}
with $\omega^\mu$ the plenum vorticity pseudovector.

In the continuum limit, the coarse-grained free energy functional becomes
\begin{equation}
\mathcal{F}[\theta]
=
\int d^4x \, dn_4
\left[
\frac{K}{2} (\partial_\mu \theta)^2
+
\frac{K_4}{2} (\partial_{n_4} \theta)^2
-
\lambda \, \Xi(x) \cos \theta
\right],
\end{equation}
where $K$ and $K_4$ are stiffness parameters along spacetime and axial directions respectively.

The axial projection of the RSVP connection is now identified as
\begin{equation}
\kappa_A(\eta)
=
\kappa_0 \, \partial_\eta \theta(\eta),
\end{equation}
so that the holonomy becomes
\begin{equation}
\beta(\eta)
=
\kappa_0 \theta(\eta_0)
-
\kappa_0 \theta(\eta).
\end{equation}

This recovers the axion-like boundary structure in the special case where $\theta$ is a propagating pseudoscalar. Here, however, $\theta$ is an emergent synchronization phase of the plenum’s parity-odd sector.

The essential quantity governing washout is the two-point correlation function of $\theta$ across the recombination kernel:
\begin{equation}
C_\theta(\Delta\eta)
=
\langle
e^{i(\theta(\eta+\Delta\eta) - \theta(\eta))}
\rangle.
\end{equation}
If the axial phase exhibits exponential decay
\begin{equation}
C_\theta(\Delta\eta)
\sim
e^{-|\Delta\eta|/\xi_{\mathrm{rec}}},
\end{equation}
then the visibility-weighted characteristic function governing polarization suppression is
\begin{equation}
\mathcal{W}_\pm
=
\left\langle
\exp
\left(
\pm 2 i \kappa_0
\left[
\theta(\eta_0) - \theta(\eta)
\right]
\right)
\right\rangle_W.
\end{equation}

Assuming small fluctuations about a coherent mean phase,
\begin{equation}
\theta(\eta)
=
\bar{\theta}(\eta)
+
\delta\theta(\eta),
\qquad
\langle \delta\theta \rangle = 0,
\end{equation}
one obtains to quadratic order
\begin{equation}
|\mathcal{W}_\pm|
=
\exp
\left(
-2 \kappa_0^2
\, \mathrm{Var}_W[\delta\theta]
\right).
\end{equation}

The variance over the visibility kernel is
\begin{equation}
\mathrm{Var}_W[\delta\theta]
=
\int d\eta\, d\eta'
W(\eta) W(\eta')
\,
\langle
\delta\theta(\eta)
\delta\theta(\eta')
\rangle.
\end{equation}

Using the exponential correlation model,
\begin{equation}
\langle
\delta\theta(\eta)
\delta\theta(\eta')
\rangle
\propto
e^{-|\eta-\eta'|/\xi_{\mathrm{rec}}},
\end{equation}
we obtain asymptotically
\begin{equation}
\mathrm{Var}_W[\delta\theta]
\sim
\frac{\Delta\eta_{\mathrm{rec}}}{\xi_{\mathrm{rec}}}
\quad
\text{for}
\quad
\xi_{\mathrm{rec}} \ll \Delta\eta_{\mathrm{rec}},
\end{equation}
and
\begin{equation}
\mathrm{Var}_W[\delta\theta]
\sim
\text{constant}
\quad
\text{for}
\quad
\xi_{\mathrm{rec}} \gg \Delta\eta_{\mathrm{rec}}.
\end{equation}

We are now in position to state the central result.

\subsection*{Coherence Theorem for the Parity-Odd Plenum Sector}

Let $\kappa_A(\eta) = \kappa_0 \partial_\eta \theta(\eta)$ be the axial projection of the RSVP connection along null geodesics, and let $\theta$ possess correlation length $\xi_{\mathrm{rec}}$ across the recombination visibility kernel of width $\Delta\eta_{\mathrm{rec}}$.

Then:

\textbf{(i)} If $\xi_{\mathrm{rec}} \gg \Delta\eta_{\mathrm{rec}}$, the parity-odd sector is coherent across recombination, and the observed birefringence angle satisfies
\begin{equation}
\beta_{\mathrm{obs}}
\approx
\kappa_0
\left(
\theta(\eta_0)
-
\theta(\eta_\star)
\right),
\end{equation}
with negligible washout suppression.

\textbf{(ii)} If $\xi_{\mathrm{rec}} \ll \Delta\eta_{\mathrm{rec}}$, the axial phase is incoherent across the emission kernel, and the observable polarization amplitude is suppressed according to
\begin{equation}
|\mathcal{W}_\pm|
=
\exp
\left(
-2 \kappa_0^2
\mathrm{Var}_W[\delta\theta]
\right),
\end{equation}
with mean rotation dominated by boundary fluctuations and therefore statistically suppressed.

\textbf{(iii)} The empirical washout bound
\begin{equation}
|\mathcal{W}_\pm|
\ge
1 - \epsilon_{\mathrm{obs}}
\end{equation}
implies a lower bound on the axial correlation length:
\begin{equation}
\xi_{\mathrm{rec}}
\gtrsim
\frac{2 \kappa_0^2 \Delta\eta_{\mathrm{rec}}}{\epsilon_{\mathrm{obs}}}.
\end{equation}

This inequality constitutes a geometric observational bound on the coherence of the parity-odd plenum sector at recombination. It does not specify the microphysical origin of $\theta$ but constrains its correlation structure.

\medskip

The two-mode axion relaxation mechanism is recovered in this framework by decomposing
\begin{equation}
\theta
=
\sum_{a=1}^N
\theta_a,
\end{equation}
with weak inter-mode correlations. The total variance is additive,
\begin{equation}
\mathrm{Var}_W[\delta\theta]
=
\sum_a
\mathrm{Var}_W[\delta\theta_a],
\end{equation}
while the coherent mean phase adds linearly. The familiar $\sqrt{N}$ relief in allowed rotation for fixed depolarization budget is thus a direct consequence of phase-variance additivity in a synchronized multi-modal axial sector.

\medskip

We have therefore replaced the particle-exclusion narrative with a coherence inequality in a synchronization field. Cosmic birefringence measures the axial holonomy of the plenum; washout measures the spectral phase noise of that sector across recombination; and the observational constraint bounds the correlation length of the parity-odd order parameter.

In the following section we discuss the implications of this coherence structure for photonic phase manipulation and entropy-limited computation.

\newpage
\section{Entropy--Coherence Inequality and Implications for Photonic and Thermodynamic Computing}

The coherence theorem derived above establishes that the observational washout constraint is equivalent to a lower bound on the axial correlation length of the parity-odd plenum sector during recombination. The same mathematical structure appears in any system in which computation is performed through controlled phase transport across a finite integration kernel. We now make this statement precise and derive an entropy--coherence inequality applicable both to cosmological birefringence and to engineered photonic computation.

\subsection*{Phase-Coherent Transport as Holonomy}

Consider a general system in which information is encoded in a phase variable $\theta(t)$ and logical transformation is achieved by controlled holonomy
\begin{equation}
\Theta_{\mathrm{out}}
=
\int_{t_i}^{t_f}
\Omega(t)\, dt,
\end{equation}
where $\Omega(t)$ is an effective angular velocity analogous to $\kappa_A$ in the cosmological case. Suppose the output observable integrates this phase over a finite response window characterized by kernel $K(t)$ of width $\Delta t$:
\begin{equation}
\mathcal{A}
=
\int dt \,
K(t)
\,
e^{i \Theta(t)}.
\end{equation}

Expanding around a coherent mean trajectory $\bar{\Theta}(t)$ with fluctuations $\delta\Theta(t)$,
\begin{equation}
\Theta(t)
=
\bar{\Theta}(t)
+
\delta\Theta(t),
\end{equation}
and assuming Gaussian phase noise for small fluctuations, one obtains
\begin{equation}
|\mathcal{A}|
\approx
\exp
\left(
-\frac{1}{2}
\mathrm{Var}_K[\delta\Theta]
\right).
\end{equation}

The preservation of computational amplitude therefore requires
\begin{equation}
\mathrm{Var}_K[\delta\Theta]
\lesssim
\varepsilon,
\end{equation}
where $\varepsilon$ is the acceptable amplitude loss.

This inequality is formally identical to the cosmological washout bound upon the identifications
\begin{equation}
t \leftrightarrow \eta,
\qquad
K(t) \leftrightarrow W(\eta),
\qquad
\Omega(t) \leftrightarrow \kappa_A(\eta).
\end{equation}

The only remaining question is what thermodynamic cost is required to maintain $\mathrm{Var}_K[\delta\Theta]$ below threshold.

\subsection*{Entropy Production and Correlation Length}

For a synchronization field governed by the free energy functional
\begin{equation}
\mathcal{F}[\theta]
=
\int
\left[
\frac{K}{2} (\nabla \theta)^2
-
\lambda \Xi \cos \theta
\right] d^dx,
\end{equation}
the correlation length $\xi$ is determined by the inverse curvature of the effective potential at the minimum:
\begin{equation}
\xi^{-2}
=
\frac{1}{K}
\left.
\frac{\partial^2 V_{\mathrm{eff}}}{\partial \theta^2}
\right|_{\theta=\theta_0}.
\end{equation}

Near criticality, the curvature vanishes and $\xi$ diverges. However, approaching criticality entails critical slowing down and increased susceptibility to noise. The fluctuation--dissipation relation yields
\begin{equation}
\langle (\delta\theta)^2 \rangle
\sim
\frac{k_B T}{K \xi^{d-2}},
\end{equation}
in $d$ effective dimensions.

For phase-coherent operation across a kernel of width $\Delta t$, we require $\xi \gg v \Delta t$, where $v$ is the characteristic propagation velocity of phase correlations. Maintaining such coherence in a noisy environment requires entropy export at a rate sufficient to counteract phase diffusion. If $D_\theta$ is the phase diffusion coefficient, the variance over the kernel scales as
\begin{equation}
\mathrm{Var}_K[\delta\Theta]
\sim
D_\theta \Delta t.
\end{equation}

The diffusion coefficient itself is constrained by the Einstein relation
\begin{equation}
D_\theta
=
\frac{k_B T}{\gamma},
\end{equation}
where $\gamma$ is the damping coefficient representing entropy export to the environment.

Combining these expressions yields the entropy--coherence inequality
\begin{equation}
\frac{k_B T}{\gamma}
\Delta t
\lesssim
\varepsilon.
\end{equation}

Equivalently,
\begin{equation}
\gamma
\gtrsim
\frac{k_B T \Delta t}{\varepsilon}.
\end{equation}

In words: the rate of entropy export (through damping) must exceed a threshold proportional to temperature and integration time in order to maintain coherent phase rotation within acceptable amplitude loss.

\subsection*{Cosmological and Computational Parallel}

In the cosmological context, $\Delta t$ corresponds to $\Delta\eta_{\mathrm{rec}}$, the thickness of recombination. The washout bound imposes
\begin{equation}
\mathrm{Var}_W[\delta\theta]
\lesssim
\varepsilon_{\mathrm{obs}},
\end{equation}
which is mathematically identical to the computational coherence constraint.

In engineered photonic computing, $\Delta t$ is the integration time of the optical detector or the effective path-length averaging window of the interferometric device. The same inequality determines the maximal logical depth achievable without active phase stabilization.

Thermodynamic computing, understood as computation constrained by entropy production and heat dissipation, therefore shares a structural equivalence with cosmological birefringence. Both are governed by the same inequality:

\begin{equation}
\text{Coherent Holonomy}
\;\;\Longleftrightarrow\;\;
\text{Phase Variance}
<
\text{Kernel Threshold}.
\end{equation}

Near-critical synchronization minimizes the energy curvature and maximizes correlation length, allowing large coherent holonomy for small control fields. However, it also increases susceptibility to phase diffusion. Optimal operation lies in a regime where $\xi$ is large relative to the kernel but finite, balancing susceptibility and damping.

\subsection*{Implications for Multi-Modal Architectures}

If the phase field decomposes into weakly coupled modes,
\begin{equation}
\theta = \sum_{a=1}^{N} \theta_a,
\end{equation}
with independent diffusion coefficients $D_a$, then
\begin{equation}
\mathrm{Var}_K[\delta\Theta]
=
\sum_a D_a \Delta t,
\end{equation}
while the coherent mean holonomy scales linearly with $\sum_a \bar{\Theta}_a$.

Thus, multi-modal phase architectures can increase total coherent rotation without increasing variance quadratically, precisely mirroring the two-field relaxation in cosmology. The square-root improvement in signal-to-noise ratio is a universal consequence of variance additivity under independent modes.

\medskip

We conclude that the washout constraint on cosmic birefringence is a cosmological instance of a general entropy--coherence principle governing phase-based computation. The parity-odd plenum sector acts as a synchronization resource whose correlation length must exceed the visibility kernel to produce rotation rather than depolarization. Photonic computing devices operate under the same inequality, with recombination thickness replaced by detector integration time and axial holonomy replaced by engineered phase gates.

The cosmological observation therefore encodes a constraint on the allowable phase noise of the universe at $z\approx1100$. In RSVP language, it bounds the entropy content of the parity-odd plenum sector during recombination. In computational language, it bounds the entropy cost of coherent phase manipulation over a finite averaging window.

\newpage
\section{Predictive Discriminants: Geometric Coherence versus Free Pseudoscalar Transport}

The preceding analysis demonstrates that cosmic birefringence may be interpreted either as the transport effect of a minimally coupled pseudoscalar field or as the axial holonomy of a parity-odd sector of the RSVP plenum connection. At the level of isotropic mean rotation and visibility-induced washout, the two descriptions are algebraically degenerate. Distinguishing them requires moving beyond the mean angle and amplitude suppression to higher-order and cross-correlational observables.

We derive three classes of discriminants: anisotropic birefringence correlations, multipole-dependent depolarization structure, and cross-correlations with entropy-gradient observables.

\subsection*{I. Anisotropic Birefringence Correlations}

Let the birefringence field on the sky be
\begin{equation}
\beta(\hat n)
=
\int_{\gamma(\hat n)}
\kappa_A(\eta, x(\eta))\, d\eta.
\end{equation}

In the minimally coupled axion model with homogeneous initial conditions and negligible clustering at recombination scales, the two-point function of $\beta$ is determined solely by the power spectrum of the axion field:
\begin{equation}
\langle \beta(\hat n)\beta(\hat n') \rangle
=
\int \frac{d^3k}{(2\pi)^3}
P_\phi(k)
e^{ik\cdot(\hat n - \hat n')\chi},
\end{equation}
with $\chi$ the comoving distance to last scattering. There is no necessary correlation with large-scale structure observables except through gravitational clustering of the axion field.

In RSVP, however, $\kappa_A$ is constructed from plenum fields such as
\begin{equation}
\kappa_A
=
\alpha \,
\varepsilon_{\mu\nu\rho\sigma}
k^\mu
u^\nu
\nabla^\rho S
v^\sigma,
\end{equation}
or similar parity-odd contractions. Consequently,
\begin{equation}
\beta(\hat n)
=
\alpha
\int_{\gamma(\hat n)}
\varepsilon_{\mu\nu\rho\sigma}
k^\mu u^\nu
\nabla^\rho S
v^\sigma
\, d\eta.
\end{equation}

The anisotropic birefringence power spectrum is therefore determined by mixed correlators of entropy gradients and vorticity:
\begin{equation}
C_\ell^{\beta\beta}
\sim
\int d\chi \,
W_\chi^2
\,
\langle
(\nabla S \cdot \omega)(\chi\hat n)
(\nabla S \cdot \omega)(\chi\hat n')
\rangle.
\end{equation}

Unlike the free axion case, this expression predicts nontrivial cross-correlations between $\beta(\hat n)$ and large-scale tracers of velocity shear, kinetic Sunyaev–Zel’dovich anisotropies, and void ellipticity statistics. In particular, one generically expects
\begin{equation}
C_\ell^{\beta\,\mathrm{kSZ}}
\neq 0,
\qquad
C_\ell^{\beta\,\nabla S}
\neq 0,
\end{equation}
provided the parity-odd plenum sector shares geometric sourcing with entropy gradients.

The absence of such correlations would favor minimally coupled axion transport; their presence would favor geometric axial coherence.

\subsection*{II. Multipole-Dependent Depolarization}

In the axion framework, the washout suppression factor
\begin{equation}
\mathcal{W}
=
J_0(g_{\phi\gamma}\phi_\star)
\end{equation}
is uniform across angular multipoles $\ell$ to leading order, since it arises from averaging over recombination thickness independent of angular scale.

In RSVP, the suppression factor depends on the correlation structure of $\theta(\eta,x)$ across both conformal time and transverse spatial separation. The visibility-weighted variance becomes
\begin{equation}
\mathrm{Var}_W[\delta\theta]
=
\int d\eta d\eta'
W(\eta)W(\eta')
\,
C_\theta(\eta,\eta';|\hat n - \hat n'|),
\end{equation}
where $C_\theta$ includes spatial dependence.

If the parity-odd sector possesses scale-dependent coherence lengths $\xi(\ell)$ in Fourier space, then the depolarization factor acquires multipole dependence:
\begin{equation}
|\mathcal{W}_\ell|
=
\exp
\left(
-2\kappa_0^2
\mathrm{Var}_W^{(\ell)}[\delta\theta]
\right).
\end{equation}

In particular, if
\begin{equation}
\xi(\ell)
\propto
\ell^{-\gamma},
\end{equation}
one expects measurable $\ell$-dependence in polarization amplitude suppression. Such scale dependence is absent in the simplest homogeneous axion model but arises naturally in a geometric plenum sector with scale-coupled entropy gradients.

Measurement of $\ell$-dependent depolarization therefore provides a discriminant between particle transport and geometric synchronization.

\subsection*{III. Cross-Correlation with Redshift-Space Entropy Gradient}

In the Everlasting Redshift framework, cosmological redshift is interpreted as an integral of entropy gradient along null geodesics:
\begin{equation}
z(\hat n)
\propto
\int_{\gamma(\hat n)}
\nabla_\mu S \, dx^\mu.
\end{equation}

If birefringence arises from the parity-odd contraction of $\nabla S$ with vorticity, then the cross-correlation
\begin{equation}
\langle \beta(\hat n) z(\hat n') \rangle
\end{equation}
is generically nonzero at second order in perturbations.

Explicitly, expanding to leading order,
\begin{equation}
\langle \beta z \rangle
\sim
\alpha
\int d\chi d\chi'
\langle
(\nabla S \cdot \omega)(\chi\hat n)
(\nabla S)(\chi'\hat n')
\rangle.
\end{equation}

In a minimally coupled axion model, no such correlation is expected unless the axion field directly couples to entropy gradients, which is not part of the canonical construction.

Detection of a statistically significant $\beta$–$z$ cross-spectrum would therefore favor the RSVP interpretation.

\subsection*{IV. Higher-Order Phase Statistics}

The axion model predicts Gaussian statistics for $\beta(\hat n)$ if the underlying field is Gaussian. In contrast, a synchronization-based parity-odd sector near criticality generically produces non-Gaussian phase fluctuations with higher-order connected correlators:
\begin{equation}
\langle \beta^3 \rangle_c \neq 0,
\qquad
\langle \beta^4 \rangle_c \propto \xi^{4-d}.
\end{equation}

Measurement of nonzero bispectrum or trispectrum in the birefringence field provides another discriminant.

\medskip

We summarize the structural difference. In the axion picture, birefringence is a boundary term of a free pseudoscalar field with minimal gravitational coupling. Its statistics and correlations are determined solely by the scalar power spectrum. In RSVP, birefringence is an axial holonomy of a parity-odd geometric sector sourced by entropy gradients and vorticity. Its statistics are therefore tied to large-scale structure observables and inherit their anisotropies, scale dependence, and higher-order correlations.

The two interpretations are degenerate at the level of isotropic mean rotation and simple washout bounds. They diverge at the level of cross-correlations, multipole dependence, and higher-order phase statistics.

\newpage
\section{Concluding Synthesis: Axial Holonomy as a Coherence Principle}

We have constructed a reinterpretation of cosmic birefringence in which the observable rotation of CMB polarization is not fundamentally the displacement of a propagating pseudoscalar degree of freedom, but the axial holonomy of a parity-odd sector of the Relativistic Scalar--Vector Plenum connection. In this framework, the Chern--Simons coupling becomes a special coordinate choice on a geometric bundle, and the axion field becomes a one-dimensional parametrization of a fiber direction within that bundle. The algebraic equivalence at the level of the mean rotation angle is preserved; the explanatory target is relocated.

The essential empirical fact underlying the washout constraint is that polarization transport is visibility-weighted. The observable rotation is not the raw holonomy
\begin{equation}
H_A[\gamma]
=
\int_{\gamma} \kappa_A \, d\eta,
\end{equation}
but its convolution with the recombination kernel. The suppression of polarization amplitude is therefore controlled not by the magnitude of $\kappa_A$ itself, but by the phase variance of the axial sector across the thickness of the emission foliation.

This leads to a coherence theorem: the observed birefringence bounds the correlation length of the parity-odd plenum sector at $z \approx 1100$. Modes coherent across the recombination thickness contribute to mean rotation; modes incoherent across that thickness contribute to depolarization. The observational constraint is thus a lower bound on axial synchronization length, not an exclusion of specific particle masses.

To render this structural claim concrete, we introduced a five-dimensional synchronization model in which the parity-odd sector is represented by a phase field $\theta$ with anisotropic couplings. The washout inequality emerges as a variance bound in a kernel-averaged characteristic function. Multi-modal relaxation of depolarization follows generically from variance additivity, independent of particle interpretation. The well-known $\sqrt{N}$ improvement in coherent rotation under multiple fields is revealed to be a universal consequence of independent phase modes, not a special property of axion multiplicity.

The cosmological constraint is therefore a specific instance of a general entropy--coherence inequality:
\begin{equation}
\mathrm{Var}_K[\delta\Theta] \lesssim \varepsilon,
\end{equation}
where $K$ denotes a finite averaging kernel. In cosmology, $K$ is the visibility function of recombination. In photonic computing, it is the integration window of an interferometric device. In thermodynamic computing, it is the temporal resolution over which phase is integrated to produce logical output. In all cases, coherent holonomy is limited by phase diffusion across the kernel.

This structural identity suggests a deeper unification. In RSVP, geometric ordering competes with entropy production. Coherent axial holonomy represents ordered structure in the parity-odd sector. Washout represents entropy-dominated phase diffusion. The CMB therefore encodes a constraint on the entropy content of the universe's parity-odd geometric degrees of freedom at recombination.

At the phenomenological level, the geometric interpretation and the minimally coupled pseudoscalar model remain degenerate for isotropic mean rotation and uniform washout. They diverge in their predictions for cross-correlations, multipole-dependent depolarization, and higher-order statistics. The RSVP interpretation generically predicts correlations between birefringence and entropy-gradient observables, as well as potential scale dependence inherited from the synchronization structure of the plenum. The axion model does not.

The conceptual shift proposed here is therefore precise. The question is no longer ``which particle produces $\beta$?'' but ``what geometric property of the plenum renders the parity-odd sector sufficiently coherent across recombination to generate the observed holonomy?'' This reframing aligns cosmic birefringence with a broader theory in which redshift, structure formation, and polarization transport are unified as manifestations of entropy-driven geometric ordering.

In this perspective, cosmic birefringence is not an isolated anomaly but a diagnostic of axial coherence in the universe’s fundamental medium. The washout bound is not merely a constraint on coupling constants; it is a statement about synchronization length in a parity-odd geometric sector. The same mathematics governs phase-stable computation in engineered photonic systems. The universe at recombination, like a photonic processor, operated within a coherence budget. The CMB polarization records that budget as a residual holonomy.

The remaining task is empirical. Should future observations reveal anisotropic correlations, scale-dependent depolarization, or non-Gaussian phase statistics in the birefringence field, the geometric coherence interpretation will gain support. Should such signatures remain absent, the minimally coupled pseudoscalar model will remain sufficient. In either case, the coherence theorem clarifies what is being tested: not merely the existence of a field, but the structure of phase synchronization in the parity-odd sector of the cosmos.

\newpage
\section*{Appendices}

\appendix

\section{Optical Transport with Axial Connection Component}

Let $k^\mu$ be a null vector satisfying
\begin{equation}
k^\mu k_\mu = 0,
\qquad
k^\nu \nabla_\nu k^\mu = 0.
\end{equation}

Let $e^\mu_a$, $a=1,2$, be orthonormal polarization vectors:
\begin{equation}
e^\mu_a k_\mu = 0,
\qquad
e^\mu_a e_{b\mu} = \delta_{ab}.
\end{equation}

Define an effective connection
\begin{equation}
\tilde{\nabla}_\nu e^\mu_a
=
\nabla_\nu e^\mu_a
+
\mathcal{A}_\nu
\epsilon_{ab}
e^\mu_b,
\end{equation}
with axial component
\begin{equation}
\kappa_A
=
\mathcal{A}_\mu k^\mu.
\end{equation}

Parallel transport yields
\begin{equation}
k^\nu \tilde{\nabla}_\nu e^\mu_a = 0
\quad \Longrightarrow \quad
\frac{d\psi}{d\eta} = \kappa_A(\eta),
\end{equation}
where $\psi$ is the polarization rotation angle.

Hence
\begin{equation}
\beta
=
\int_{\gamma}
\kappa_A(\eta)\, d\eta.
\end{equation}

\section{Visibility-Weighted Holonomy}

Let $W(\eta)$ be normalized:
\begin{equation}
\int W(\eta)\, d\eta = 1.
\end{equation}

Define
\begin{equation}
\beta(\eta)
=
\int_{\eta}^{\eta_0}
\kappa_A(\eta')\, d\eta'.
\end{equation}

Observable rotation:
\begin{equation}
\beta_{\mathrm{obs}}
=
\int W(\eta)\beta(\eta)\, d\eta.
\end{equation}

Interchanging integrals,
\begin{equation}
\beta_{\mathrm{obs}}
=
\int \kappa_A(\eta')
\left(
\int^{\eta'} W(\eta)\, d\eta
\right)
d\eta'.
\end{equation}

\section{Phase Variance and Depolarization}

Define
\begin{equation}
\Theta(\eta)
=
2\beta(\eta).
\end{equation}

Characteristic function:
\begin{equation}
\mathcal{W}_\pm
=
\left\langle
e^{\pm i \Theta}
\right\rangle_W.
\end{equation}

Assuming Gaussian fluctuations:
\begin{equation}
|\mathcal{W}_\pm|
=
\exp
\left(
-\frac{1}{2}
\mathrm{Var}_W[\Theta]
\right).
\end{equation}

Variance:
\begin{equation}
\mathrm{Var}_W[\Theta]
=
4
\int d\eta d\eta'
W(\eta)W(\eta')
\,
\mathrm{Cov}\!\left(
\beta(\eta),\beta(\eta')
\right).
\end{equation}

For stationary fluctuations:
\begin{equation}
\mathrm{Cov}(\beta(\eta),\beta(\eta'))
=
\int d\tau d\tau'
\langle
\kappa_A(\tau)\kappa_A(\tau')
\rangle.
\end{equation}

If
\begin{equation}
\langle
\kappa_A(\eta)\kappa_A(\eta')
\rangle
\propto
e^{-|\eta-\eta'|/\xi},
\end{equation}
then
\begin{equation}
\mathrm{Var}_W[\Theta]
\sim
\frac{\Delta\eta_{\mathrm{rec}}}{\xi}.
\end{equation}

\section{Multi-Modal Decomposition}

Let
\begin{equation}
\kappa_A
=
\sum_{a=1}^{N}
\kappa_A^{(a)},
\qquad
\langle \kappa_A^{(a)} \kappa_A^{(b)} \rangle
=
0
\quad (a\neq b).
\end{equation}

Then
\begin{equation}
\mathrm{Var}_W[\Theta]
=
\sum_{a=1}^{N}
\mathrm{Var}_W[\Theta_a].
\end{equation}

Mean rotation:
\begin{equation}
\langle \beta \rangle
=
\sum_{a=1}^{N}
\langle \beta_a \rangle.
\end{equation}

Signal-to-variance scaling:
\begin{equation}
\frac{\langle \beta \rangle}{\sqrt{\mathrm{Var}_W}}
\propto
\sqrt{N}.
\end{equation}

\section{Synchronization Field Reduction}

Let $\theta$ satisfy
\begin{equation}
\mathcal{F}[\theta]
=
\int
\left[
\frac{K}{2}(\nabla\theta)^2
-
\lambda \Xi \cos\theta
\right] d^dx.
\end{equation}

Define
\begin{equation}
\kappa_A
=
\kappa_0 \partial_\eta \theta.
\end{equation}

Correlation function:
\begin{equation}
\langle
e^{i(\theta(x)-\theta(y))}
\rangle
=
e^{-|x-y|/\xi}.
\end{equation}

Then
\begin{equation}
\mathrm{Var}_W[\delta\theta]
=
\int d\eta d\eta'
W(\eta)W(\eta')
e^{-|\eta-\eta'|/\xi}.
\end{equation}

\section{Entropy--Coherence Inequality}

Let phase diffusion obey
\begin{equation}
\langle (\delta\theta)^2 \rangle
=
2D_\theta t.
\end{equation}

Fluctuation--dissipation:
\begin{equation}
D_\theta
=
\frac{k_B T}{\gamma}.
\end{equation}

Coherence condition:
\begin{equation}
\mathrm{Var}_K[\delta\Theta]
=
2D_\theta \Delta t
\le \varepsilon.
\end{equation}

Hence
\begin{equation}
\gamma
\ge
\frac{2k_B T \Delta t}{\varepsilon}.
\end{equation}

\section{Axial Coherence Bound at Recombination}

Observed suppression:
\begin{equation}
|\mathcal{W}_\pm|
\ge
1-\epsilon_{\mathrm{obs}}.
\end{equation}

Thus
\begin{equation}
\mathrm{Var}_W[\Theta]
\le
2\epsilon_{\mathrm{obs}}.
\end{equation}

For exponential correlations:
\begin{equation}
\xi_{\mathrm{rec}}
\gtrsim
\frac{\Delta\eta_{\mathrm{rec}}}{2\epsilon_{\mathrm{obs}}}.
\end{equation}

\section{Spectral Formulation of Visibility-Weighted Axial Holonomy}

Let conformal time be denoted $\eta$ and define the axial projection
\begin{equation}
\kappa_A(\eta,\mathbf{x})
=
\kappa_0 \, \partial_\eta \theta(\eta,\mathbf{x}).
\end{equation}

Define Fourier transform in comoving coordinates:
\begin{equation}
\theta(\eta,\mathbf{x})
=
\int \frac{d^3k}{(2\pi)^3}
\, \theta_{\mathbf{k}}(\eta)
\, e^{i \mathbf{k}\cdot \mathbf{x}}.
\end{equation}

The axial power spectrum is
\begin{equation}
\langle
\theta_{\mathbf{k}}(\eta)
\theta_{\mathbf{k}'}^*(\eta')
\rangle
=
(2\pi)^3
\delta^{(3)}(\mathbf{k}-\mathbf{k}')
P_\theta(k;\eta,\eta').
\end{equation}

Then
\begin{equation}
\langle
\kappa_A(\eta)
\kappa_A(\eta')
\rangle
=
\kappa_0^2
\partial_\eta \partial_{\eta'}
\int \frac{d^3k}{(2\pi)^3}
P_\theta(k;\eta,\eta').
\end{equation}

Define the line-of-sight holonomy
\begin{equation}
\beta(\hat n)
=
\int d\eta \,
\kappa_A(\eta,\chi\hat n).
\end{equation}

Angular power spectrum:
\begin{equation}
C_\ell^{\beta\beta}
=
\int \frac{2}{\pi} k^2 dk
\left|
\int d\eta \,
j_\ell(k\chi)
\, \kappa_0 \partial_\eta T_\theta(k,\eta)
\right|^2
P_\theta^{\mathrm{prim}}(k),
\end{equation}
where $T_\theta$ is the transfer function.

Visibility weighting modifies the kernel:
\begin{equation}
\beta_{\mathrm{obs}}
=
\int d\eta \,
W(\eta)
\int_{\eta}^{\eta_0}
\kappa_A(\eta') d\eta'.
\end{equation}

Fourier-space variance:
\begin{equation}
\mathrm{Var}_W[\Theta]
=
4\kappa_0^2
\int \frac{d^3k}{(2\pi)^3}
\left|
\int d\eta \,
W(\eta)
\int_{\eta}^{\eta_0}
\partial_{\eta'} T_\theta(k,\eta')
d\eta'
\right|^2
P_\theta^{\mathrm{prim}}(k).
\end{equation}

Define spectral coherence scale:
\begin{equation}
k_{\mathrm{coh}}
\sim
\frac{1}{\Delta\eta_{\mathrm{rec}}}.
\end{equation}

Modes with $k \gg k_{\mathrm{coh}}$ are suppressed by oscillatory cancellation in the double integral; modes with $k \ll k_{\mathrm{coh}}$ contribute coherently.

Hence the washout bound constrains
\begin{equation}
\int_{k>k_{\mathrm{coh}}}
P_\theta^{\mathrm{prim}}(k)
\, dk
\le
\mathcal{O}(\epsilon_{\mathrm{obs}}).
\end{equation}

\section{Renormalization-Group Analysis of Axial Phase Synchronization}

Consider the Euclidean action for the axial phase field in $d$ spatial dimensions:

\begin{equation}
S[\theta]
=
\int d^d x
\left[
\frac{K}{2} (\nabla \theta)^2
-
\lambda \, \Xi(x) \cos \theta
\right].
\end{equation}

For $\Xi=0$ this reduces to the $O(2)$ (XY) model. Introduce dimensionless couplings:

\begin{equation}
g = \frac{T}{K},
\qquad
y = \frac{\lambda}{T}.
\end{equation}

Under coarse-graining $x \rightarrow x' = x/b$ with $b>1$, rescale field $\theta'(x')=\theta(x)$.

Momentum-shell integration over $\Lambda/b < |\mathbf{k}| < \Lambda$ yields one-loop beta functions (for $d=4-\epsilon$):

\begin{equation}
\beta_g
=
\frac{dg}{d\ln b}
=
-\epsilon g
+
C_d g^2
+
\mathcal{O}(g^3),
\end{equation}

\begin{equation}
\beta_y
=
\frac{dy}{d\ln b}
=
\left(
2-\frac{g}{4\pi^2}
\right)y
+
\mathcal{O}(gy^2).
\end{equation}

Fixed points satisfy

\begin{equation}
\beta_g=0,
\qquad
\beta_y=0.
\end{equation}

Gaussian fixed point:
\begin{equation}
g^\ast=0,
\qquad
y^\ast=0.
\end{equation}

Wilson-Fisher fixed point:
\begin{equation}
g^\ast=\frac{\epsilon}{C_d},
\qquad
y^\ast=0.
\end{equation}

Correlation length near criticality:

\begin{equation}
\xi
\sim
|g-g^\ast|^{-\nu},
\qquad
\nu
=
\frac{1}{2}
+
\mathcal{O}(\epsilon).
\end{equation}

Axial coherence across recombination requires

\begin{equation}
\xi_{\mathrm{rec}}
\gg
\Delta\eta_{\mathrm{rec}}.
\end{equation}

Define reduced temperature parameter:

\begin{equation}
t
=
\frac{g-g^\ast}{g^\ast}.
\end{equation}

Then

\begin{equation}
\xi_{\mathrm{rec}}
\sim
|t|^{-\nu}.
\end{equation}

Washout bound implies

\begin{equation}
|t|
\lesssim
\left(
\frac{\Delta\eta_{\mathrm{rec}}}{L_0}
\right)^{-1/\nu},
\end{equation}

where $L_0$ is microscopic scale.

Phase stiffness renormalization:

\begin{equation}
K_R(b)
=
K b^{-\eta},
\qquad
\eta
=
\mathcal{O}(\epsilon^2).
\end{equation}

Phase variance scaling:

\begin{equation}
\langle (\delta\theta)^2 \rangle
\sim
\int^{\Lambda}
\frac{d^d k}{(2\pi)^d}
\frac{T}{K_R k^2}
\sim
\frac{T}{K_R}
\Lambda^{d-2}.
\end{equation}

Near criticality:

\begin{equation}
\langle (\delta\theta)^2 \rangle
\propto
\xi^{2-d}.
\end{equation}

For effective $d<4$, coherence enhances as $d \rightarrow 2$.

Multi-modal decomposition:

\begin{equation}
\theta
=
\sum_{a=1}^N \theta_a,
\end{equation}

with independent RG flows:

\begin{equation}
\beta_{g_a}
=
-\epsilon g_a + C_d g_a^2.
\end{equation}

Total variance:

\begin{equation}
\mathrm{Var}
=
\sum_a \mathrm{Var}_a.
\end{equation}

Signal-to-noise scaling:

\begin{equation}
\frac{\langle \beta \rangle}{\sqrt{\mathrm{Var}}}
\sim
\sqrt{N}
|t|^{\nu}.
\end{equation}

\section{Principal Bundle Formulation of the Axial Connection Sector}

Let $M$ be a four-dimensional Lorentzian manifold with metric $g_{\mu\nu}$.
Let $P(M,G)$ be a principal $G$-bundle over $M$, with structure group
\begin{equation}
G = U(1)_A,
\end{equation}
representing the axial phase sector.

Let $\mathcal{A}$ denote a connection one-form on $P$:
\begin{equation}
\mathcal{A} \in \Omega^1(P,\mathfrak{u}(1)).
\end{equation}

Pull back to spacetime via a local section $s: M \to P$:
\begin{equation}
A_\mu = s^\ast \mathcal{A}.
\end{equation}

Curvature two-form:
\begin{equation}
\mathcal{F} = d\mathcal{A},
\qquad
F_{\mu\nu} = \partial_\mu A_\nu - \partial_\nu A_\mu.
\end{equation}

Axial projection along null direction $k^\mu$:

\begin{equation}
\kappa_A = A_\mu k^\mu.
\end{equation}

Holonomy along null geodesic $\gamma$:

\begin{equation}
\mathrm{Hol}_\gamma(A)
=
\exp
\left(
i \int_\gamma A
\right).
\end{equation}

Polarization rotation:

\begin{equation}
\beta = \int_\gamma A_\mu k^\mu d\eta.
\end{equation}

Gauge transformation:

\begin{equation}
A_\mu \to A_\mu + \partial_\mu \alpha,
\qquad
\theta \to \theta + \alpha.
\end{equation}

Holonomy gauge invariance:

\begin{equation}
\int_\gamma \partial_\mu \alpha \, dx^\mu
=
\alpha(\eta_0) - \alpha(\eta_\star).
\end{equation}

Visibility-weighted observable:

\begin{equation}
\beta_{\mathrm{obs}}
=
\int d\eta \,
W(\eta)
\int_{\eta}^{\eta_0}
A_\mu k^\mu d\eta'.
\end{equation}

Bundle curvature interpretation:

\begin{equation}
\int_\Sigma F
=
\int_{\partial \Sigma} A,
\end{equation}

where $\Sigma$ is a 2-surface bounded by $\gamma$ and reference curve.

Parity-odd construction from plenum fields:

\begin{equation}
A_\mu
=
\alpha
\varepsilon_{\mu\nu\rho\sigma}
u^\nu
\nabla^\rho S
v^\sigma.
\end{equation}

Then curvature:

\begin{equation}
F_{\mu\nu}
=
\alpha
\varepsilon_{\mu\nu\rho\sigma}
\left(
\nabla^\rho \nabla^\lambda S
\right)
v_\lambda
+
\text{antisymmetric terms}.
\end{equation}

Derived categorical formulation:

Let $\mathbf{Conn}(M)$ denote the groupoid of connections on $P$.
Objects: $(P,\mathcal{A})$.
Morphisms: gauge transformations.

Define holonomy functor:

\begin{equation}
\mathrm{Hol}:
\Pi_1(M)
\to
U(1)_A,
\end{equation}

where $\Pi_1(M)$ is the fundamental groupoid.

Visibility weighting defines natural transformation:

\begin{equation}
\mathcal{V}:
\mathrm{Hol}
\Rightarrow
\mathrm{Hol}_W,
\end{equation}

with

\begin{equation}
\mathrm{Hol}_W(\gamma)
=
\exp
\left(
i \int W(\eta) A_\mu k^\mu d\eta
\right).
\end{equation}

Coherence condition in bundle form:

\begin{equation}
\left|
\left\langle
\mathrm{Hol}_W(\gamma)
\right\rangle
\right|
\ge
1-\epsilon.
\end{equation}

Phase variance expressed via curvature two-point function:

\begin{equation}
\mathrm{Var}(\beta)
=
\int d\eta d\eta'
W(\eta)W(\eta')
\langle
A_\mu(\eta) A_\nu(\eta')
\rangle
k^\mu k^\nu.
\end{equation}

Correlation length in bundle language:

\begin{equation}
\langle
A_\mu(x) A_\nu(y)
\rangle
\sim
e^{-|x-y|/\xi}.
\end{equation}

Coherence bound:

\begin{equation}
\xi
\gg
\Delta\eta_{\mathrm{rec}}.
\end{equation}

\section{Derived-Stack Formulation of the Axial Connection Sector}

Let $M$ be a smooth Lorentzian manifold.
Define the moduli stack of $U(1)_A$-connections:

\begin{equation}
\mathcal{Conn}_{U(1)_A}(M)
=
\left[
\Omega^1(M)
\Big/
\Omega^0(M)
\right],
\end{equation}

where $\Omega^0(M)$ acts by gauge transformation:

\begin{equation}
A \mapsto A + d\alpha.
\end{equation}

Regard $\mathcal{Conn}_{U(1)_A}(M)$ as a derived stack via the mapping stack:

\begin{equation}
\mathbb{R}\mathrm{Map}(M, B U(1)_A).
\end{equation}

Tangent complex at a connection $A$:

\begin{equation}
T_A
\mathcal{Conn}
\simeq
\left[
\Omega^0(M)
\xrightarrow{d}
\Omega^1(M)
\right],
\end{equation}

placed in degrees $[-1,0]$.

Derived curvature map:

\begin{equation}
\mathcal{F}:
\mathcal{Conn}
\to
\Omega^2(M),
\qquad
A \mapsto dA.
\end{equation}

Null geodesic $\gamma$ defines morphism in the fundamental $\infty$-groupoid:

\begin{equation}
\gamma: \Delta^1 \to M.
\end{equation}

Holonomy as derived pullback:

\begin{equation}
\mathrm{Hol}_\gamma(A)
=
\exp
\left(
i \int_{\Delta^1}
\gamma^\ast A
\right).
\end{equation}

Visibility kernel defines weighted integration functional:

\begin{equation}
\mathcal{I}_W:
\Omega^1(M)
\to
\mathbb{R},
\qquad
A \mapsto
\int W(\eta)\, \gamma^\ast A.
\end{equation}

Thus observable birefringence:

\begin{equation}
\beta_{\mathrm{obs}}
=
\mathcal{I}_W(A).
\end{equation}

Derived deformation complex for axial phase field $\theta$:

\begin{equation}
\mathbb{T}_\theta
\simeq
\left[
C^\infty(M)
\xrightarrow{\Box + V''(\theta)}
C^\infty(M)
\right].
\end{equation}

Correlation length from derived propagator:

\begin{equation}
G(x,y)
=
\left(
\Box + m_{\mathrm{eff}}^2
\right)^{-1},
\qquad
\xi = m_{\mathrm{eff}}^{-1}.
\end{equation}

Homotopy-coherent multi-modal decomposition:

\begin{equation}
A
=
\bigoplus_{a=1}^N
A^{(a)}
\quad
\text{in}
\quad
\mathrm{Perf}(\mathcal{Conn}),
\end{equation}

with derived direct sum preserving variance additivity:

\begin{equation}
\mathrm{Var}(\beta)
=
\sum_a
\mathrm{Var}_a(\beta).
\end{equation}

Coherence condition as derived inequality:

\begin{equation}
\| \mathcal{I}_W(A) \|_{L^2}
\le
\epsilon
\quad
\Longrightarrow
\quad
\xi
\gg
\Delta\eta_{\mathrm{rec}}.
\end{equation}

Homotopy fiber of visibility functional:

\begin{equation}
\mathrm{hofib}(\mathcal{I}_W)
=
\left\{
A \in \mathcal{Conn}
\mid
\mathcal{I}_W(A)=0
\right\}.
\end{equation}

Derived enhancement of coherence bound:

\begin{equation}
\mathrm{Spec}
\left(
\mathbb{T}_\theta
\right)
\cap
\left[
k > \frac{1}{\Delta\eta_{\mathrm{rec}}}
\right]
=
\varnothing
\quad
\text{(coherent sector)}.
\end{equation}

Functorial statement:

\begin{equation}
\mathcal{V}:
\mathbb{R}\mathrm{Map}(M,B U(1)_A)
\to
\mathbb{R},
\end{equation}

with coherence constraint:

\begin{equation}
\mathrm{Im}(\mathcal{V})
\subset
[-\epsilon,\epsilon].
\end{equation}

\section{Stochastic Differential Equation Formulation of Axial Phase Dynamics}

Let the axial phase field $\theta(\eta,\mathbf{x})$ satisfy a Langevin-type stochastic evolution equation:

\begin{equation}
\partial_\eta \theta
=
- \Gamma \frac{\delta \mathcal{F}}{\delta \theta}
+
\zeta(\eta,\mathbf{x}),
\end{equation}

where $\Gamma$ is a damping coefficient and $\zeta$ is Gaussian white noise with

\begin{equation}
\langle \zeta(\eta,\mathbf{x}) \rangle = 0,
\end{equation}

\begin{equation}
\langle
\zeta(\eta,\mathbf{x})
\zeta(\eta',\mathbf{x}')
\rangle
=
2\Gamma T
\delta(\eta-\eta')
\delta^{(3)}(\mathbf{x}-\mathbf{x}').
\end{equation}

Free energy functional:

\begin{equation}
\mathcal{F}[\theta]
=
\int d^3x
\left[
\frac{K}{2} (\nabla \theta)^2
+
\frac{m^2}{2} \theta^2
\right].
\end{equation}

Linearized dynamics:

\begin{equation}
\partial_\eta \theta
=
- \Gamma ( -K\nabla^2 + m^2 ) \theta
+
\zeta.
\end{equation}

Fourier decomposition:

\begin{equation}
\theta(\eta,\mathbf{x})
=
\int \frac{d^3k}{(2\pi)^3}
\theta_{\mathbf{k}}(\eta)
e^{i\mathbf{k}\cdot\mathbf{x}}.
\end{equation}

Mode evolution:

\begin{equation}
\partial_\eta \theta_{\mathbf{k}}
=
- \Gamma (K k^2 + m^2)\theta_{\mathbf{k}}
+
\zeta_{\mathbf{k}}.
\end{equation}

Solution:

\begin{equation}
\theta_{\mathbf{k}}(\eta)
=
\theta_{\mathbf{k}}(0)
e^{-\Gamma (Kk^2+m^2)\eta}
+
\int_0^\eta
e^{-\Gamma (Kk^2+m^2)(\eta-s)}
\zeta_{\mathbf{k}}(s)
ds.
\end{equation}

Two-point function:

\begin{equation}
\langle
\theta_{\mathbf{k}}(\eta)
\theta_{\mathbf{k}'}(\eta')
\rangle
=
(2\pi)^3 \delta^{(3)}(\mathbf{k}+\mathbf{k}')
\frac{T}{Kk^2+m^2}
e^{-\Gamma (Kk^2+m^2)|\eta-\eta'|}.
\end{equation}

Correlation length:

\begin{equation}
\xi = m^{-1}.
\end{equation}

Phase diffusion coefficient:

\begin{equation}
D_\theta(k)
=
\frac{T}{\Gamma (Kk^2+m^2)}.
\end{equation}

Holonomy:

\begin{equation}
\beta(\eta)
=
\kappa_0 \int_\eta^{\eta_0}
\partial_{\eta'} \theta(\eta') d\eta'.
\end{equation}

Variance:

\begin{equation}
\mathrm{Var}_W[\Theta]
=
4\kappa_0^2
\int d\eta d\eta'
W(\eta)W(\eta')
\langle
\delta\theta(\eta)
\delta\theta(\eta')
\rangle.
\end{equation}

Using exponential temporal decay:

\begin{equation}
\langle
\delta\theta(\eta)
\delta\theta(\eta')
\rangle
\sim
e^{-\Gamma m^2 |\eta-\eta'|}.
\end{equation}

Effective coherence time:

\begin{equation}
\tau_c
=
(\Gamma m^2)^{-1}.
\end{equation}

Visibility variance scaling:

\begin{equation}
\mathrm{Var}_W[\Theta]
\sim
\kappa_0^2
\frac{\Delta\eta_{\mathrm{rec}}}{\tau_c}.
\end{equation}

Coherence condition:

\begin{equation}
\tau_c
\gg
\Delta\eta_{\mathrm{rec}}.
\end{equation}

Equivalent inequality:

\begin{equation}
\Gamma m^2
\ll
\frac{1}{\Delta\eta_{\mathrm{rec}}}.
\end{equation}

Entropy production rate:

\begin{equation}
\dot{S}
=
\int d^3x
\frac{1}{T}
\left(
\partial_\eta \theta
\right)^2.
\end{equation}

Stationary regime:

\begin{equation}
\dot{S}
\propto
\Gamma
\langle (\delta\theta)^2 \rangle.
\end{equation}

Entropy--coherence relation:

\begin{equation}
\dot{S}
\lesssim
\frac{T}{\Delta\eta_{\mathrm{rec}}}.
\end{equation}

\section{Chern--Simons Effective Action and Axial Projection}

Consider the electromagnetic field strength
\begin{equation}
F_{\mu\nu} = \partial_\mu A_\nu - \partial_\nu A_\mu.
\end{equation}

Introduce pseudoscalar field $\phi$ with Chern--Simons coupling:

\begin{equation}
S_{\mathrm{CS}}
=
\int d^4x
\left[
-\frac{1}{4} F_{\mu\nu}F^{\mu\nu}
-
\frac{g}{4}
\phi
F_{\mu\nu}\tilde{F}^{\mu\nu}
\right],
\end{equation}

where
\begin{equation}
\tilde{F}^{\mu\nu}
=
\frac{1}{2}
\varepsilon^{\mu\nu\rho\sigma}
F_{\rho\sigma}.
\end{equation}

Modified Maxwell equations:

\begin{equation}
\nabla_\mu F^{\mu\nu}
=
g
(\partial_\mu \phi)
\tilde{F}^{\mu\nu}.
\end{equation}

Geometric optics ansatz:

\begin{equation}
A_\mu = \Re \left[ a_\mu e^{i S/\epsilon} \right],
\end{equation}

with wavevector
\begin{equation}
k_\mu = \partial_\mu S.
\end{equation}

Leading-order polarization transport:

\begin{equation}
k^\nu \nabla_\nu a_\mu
=
\frac{g}{2}
(\partial_\nu \phi)
\varepsilon^{\nu}{}_{\mu\rho\sigma}
k^\rho a^\sigma.
\end{equation}

Define effective axial connection:

\begin{equation}
A_\mu^{(A)}
=
\frac{g}{2}
\partial_\mu \phi.
\end{equation}

Then

\begin{equation}
\kappa_A
=
A_\mu^{(A)} k^\mu.
\end{equation}

Holonomy:

\begin{equation}
\beta
=
\frac{g}{2}
\left(
\phi(\eta_0)
-
\phi(\eta_\star)
\right).
\end{equation}

Generalized parity-odd plenum connection:

\begin{equation}
A_\mu^{(A)}
=
\alpha
\varepsilon_{\mu\nu\rho\sigma}
u^\nu
\nabla^\rho S
v^\sigma.
\end{equation}

Effective Chern--Simons term:

\begin{equation}
S_{\mathrm{eff}}
=
\int
A_\mu^{(A)} J^\mu_{\mathrm{EM}}
d^4x,
\end{equation}

with
\begin{equation}
J^\mu_{\mathrm{EM}}
=
\varepsilon^{\mu\nu\rho\sigma}
A_\nu F_{\rho\sigma}.
\end{equation}

\section{Cumulant Expansion of Axial Phase Fluctuations}

Define phase:

\begin{equation}
\Theta
=
2\beta.
\end{equation}

Characteristic function:

\begin{equation}
\mathcal{W}
=
\langle e^{i\Theta} \rangle.
\end{equation}

Cumulant expansion:

\begin{equation}
\ln \mathcal{W}
=
\sum_{n=1}^{\infty}
\frac{(i)^n}{n!}
\kappa_n,
\end{equation}

where $\kappa_n$ are cumulants of $\Theta$.

First cumulant:

\begin{equation}
\kappa_1 = \langle \Theta \rangle.
\end{equation}

Second cumulant:

\begin{equation}
\kappa_2 = \langle \Theta^2 \rangle - \langle \Theta \rangle^2.
\end{equation}

Third cumulant:

\begin{equation}
\kappa_3
=
\langle (\Theta-\langle\Theta\rangle)^3 \rangle.
\end{equation}

Depolarization factor including skewness:

\begin{equation}
|\mathcal{W}|
=
\exp
\left(
-\frac{\kappa_2}{2}
+
\frac{\kappa_4}{24}
-
\cdots
\right).
\end{equation}

If axial sector near criticality:

\begin{equation}
\kappa_n
\propto
\xi^{(n-1)(2-d)}.
\end{equation}

Bispectrum contribution to birefringence:

\begin{equation}
\langle
\beta(\hat n_1)
\beta(\hat n_2)
\beta(\hat n_3)
\rangle
=
\int
\prod_i d\eta_i
W(\eta_i)
\langle
\kappa_A(\eta_1)
\kappa_A(\eta_2)
\kappa_A(\eta_3)
\rangle.
\end{equation}

Non-zero $\kappa_3$ implies parity-odd non-Gaussianity.

\section{Topological and Torsion-Class Representation of the Axial Sector}

Let $M$ admit orthonormal frame $e^a$ with spin connection $\omega^{ab}$.

Torsion two-form:

\begin{equation}
T^a
=
de^a
+
\omega^a{}_b \wedge e^b.
\end{equation}

Parity-odd invariant:

\begin{equation}
\mathcal{P}
=
\varepsilon_{abcd}
T^a \wedge T^b \wedge e^c \wedge e^d.
\end{equation}

Define axial one-form:

\begin{equation}
A^{(A)}
=
\ast (T^a \wedge e_a).
\end{equation}

Holonomy:

\begin{equation}
\beta
=
\int_\gamma A^{(A)}.
\end{equation}

Chern class:

\begin{equation}
c_1(P)
=
\frac{1}{2\pi}
\int F.
\end{equation}

Parity-odd topological density:

\begin{equation}
\mathcal{L}_{\mathrm{CS}}
=
\phi
F \wedge F.
\end{equation}

Generalized plenum invariant:

\begin{equation}
\mathcal{I}
=
\int_M
\varepsilon^{\mu\nu\rho\sigma}
\nabla_\mu S
\nabla_\nu v_\rho
u_\sigma
d^4x.
\end{equation}

Coherence constraint:

\begin{equation}
\left|
\int_\gamma A^{(A)}
\right|
\le
\epsilon
\quad
\Longrightarrow
\quad
[T^a] \text{ torsion class bounded}.
\end{equation}

\section{Spectral Determinants, Stochastic Path Integrals, and Categorical Entropy Flow}

\subsection{Spectral Zeta-Function and Axial Fluctuation Determinant}

Let the quadratic fluctuation operator of the axial phase be

\begin{equation}
\mathcal{O}
=
-\nabla^2 + m^2.
\end{equation}

Eigenvalue problem:

\begin{equation}
\mathcal{O} \psi_n = \lambda_n \psi_n.
\end{equation}

Define spectral zeta function:

\begin{equation}
\zeta_{\mathcal{O}}(s)
=
\sum_n \lambda_n^{-s}.
\end{equation}

Regularized determinant:

\begin{equation}
\log \det \mathcal{O}
=
- \left. \frac{d}{ds} \zeta_{\mathcal{O}}(s) \right|_{s=0}.
\end{equation}

One-loop effective action:

\begin{equation}
\Gamma_{\mathrm{eff}}
=
\frac{1}{2}
\log \det \mathcal{O}.
\end{equation}

Heat kernel representation:

\begin{equation}
\zeta_{\mathcal{O}}(s)
=
\frac{1}{\Gamma(s)}
\int_0^\infty
t^{s-1}
\mathrm{Tr}
\left(
e^{-t\mathcal{O}}
\right)
dt.
\end{equation}

Short-time expansion:

\begin{equation}
\mathrm{Tr}
\left(
e^{-t\mathcal{O}}
\right)
=
\frac{1}{(4\pi t)^{d/2}}
\sum_{k=0}^\infty
a_k t^k.
\end{equation}

Correlation length:

\begin{equation}
\xi = m^{-1}.
\end{equation}

High-frequency modes ($\lambda_n \gg \xi^{-2}$) suppressed in effective action.

Coherence condition:

\begin{equation}
\sum_{\lambda_n > \Delta\eta_{\mathrm{rec}}^{-2}}
\lambda_n^{-1}
\le
\epsilon.
\end{equation}

\subsection{Martin--Siggia--Rose Functional for Axial Phase Diffusion}

Stochastic equation:

\begin{equation}
\partial_\eta \theta
=
- \Gamma \frac{\delta \mathcal{F}}{\delta \theta}
+
\zeta.
\end{equation}

Noise:

\begin{equation}
\langle \zeta(\eta,x)\zeta(\eta',x') \rangle
=
2\Gamma T \delta(\eta-\eta')\delta(x-x').
\end{equation}

MSR generating functional:

\begin{equation}
Z
=
\int \mathcal{D}\theta \mathcal{D}\hat{\theta}
\exp
\left(
- S_{\mathrm{MSR}}[\theta,\hat{\theta}]
\right).
\end{equation}

Action:

\begin{equation}
S_{\mathrm{MSR}}
=
\int d\eta d^3x
\left[
\hat{\theta}
\left(
\partial_\eta \theta
+
\Gamma \frac{\delta \mathcal{F}}{\delta \theta}
\right)
-
\Gamma T \hat{\theta}^2
\right].
\end{equation}

Two-point response function:

\begin{equation}
G_R(k,\omega)
=
\frac{1}{-i\omega + \Gamma (Kk^2+m^2)}.
\end{equation}

Phase variance:

\begin{equation}
\langle |\theta_k|^2 \rangle
=
\frac{T}{Kk^2+m^2}.
\end{equation}

Visibility-weighted cumulant:

\begin{equation}
\mathrm{Var}_W[\Theta]
=
4\kappa_0^2
\int \frac{d^3k}{(2\pi)^3}
\frac{T}{Kk^2+m^2}
|\tilde{W}(k)|^2.
\end{equation}

Coherence inequality:

\begin{equation}
\int_{k > 1/\Delta\eta_{\mathrm{rec}}}
\frac{T}{Kk^2+m^2}
dk
\le
\epsilon.
\end{equation}

\subsection{Categorical Entropy-Flow Formalism}

Define category $\mathbf{Ax}$:

Objects: axial phase configurations $\theta$.

Morphisms: evolution maps
\begin{equation}
\Phi_{\eta_1\to\eta_2}:
\theta(\eta_1)\to\theta(\eta_2).
\end{equation}

Entropy functional:

\begin{equation}
\mathcal{S}(\theta)
=
- \int P[\theta]\log P[\theta].
\end{equation}

Define functor:

\begin{equation}
\mathcal{E}:
\mathbf{Ax}
\to
\mathbf{R}_{\ge 0},
\qquad
\theta \mapsto \dot{S}.
\end{equation}

Entropy production rate:

\begin{equation}
\dot{S}
=
\int
\frac{1}{T}
(\partial_\eta \theta)^2 d^3x.
\end{equation}

Holonomy functional:

\begin{equation}
\mathcal{H}(\theta)
=
\int_\gamma \kappa_0 \partial_\eta \theta d\eta.
\end{equation}

Define coherence morphism constraint:

\begin{equation}
|\mathcal{H}(\theta)|
\ge
\beta_{\mathrm{obs}},
\qquad
\dot{S}
\le
S_{\mathrm{max}}.
\end{equation}

Entropy--holonomy inequality:

\begin{equation}
\beta_{\mathrm{obs}}^2
\le
\kappa_0^2
\Delta\eta_{\mathrm{rec}}
\int \frac{\dot{S}}{T} d\eta.
\end{equation}

Functorial composition law:

\begin{equation}
\mathcal{H}(\theta_2\circ\theta_1)
=
\mathcal{H}(\theta_2)
+
\mathcal{H}(\theta_1).
\end{equation}

Entropy monotonicity:

\begin{equation}
\mathcal{E}(\theta_2\circ\theta_1)
\ge
\mathcal{E}(\theta_1).
\end{equation}

Coherence admissibility condition:

\begin{equation}
\exists \theta \in \mathbf{Ax}
\quad
\text{s.t.}
\quad
\mathcal{H}(\theta)=\beta_{\mathrm{obs}}
\quad
\land
\quad
\mathcal{E}(\theta)\le S_{\mathrm{max}}.
\end{equation}



\newpage
\begin{thebibliography}{99}

\bibitem{ShaoObataZhang2025}
J.~Shao, I.~Obata and D.~Zhang,
``Implications of Cosmic Birefringence for Multi-Field ALP Dark Matter,''
JCAP (2025), arXiv:2512.21888 [astro-ph.CO].

\bibitem{Carroll1990}
S.~M.~Carroll, G.~B.~Field and R.~Jackiw,
``Limits on a Lorentz- and parity-violating modification of electrodynamics,''
Phys.\ Rev.\ D {\bf 41} (1990) 1231.

\bibitem{Lue1999}
A.~Lue, L.~Wang and M.~Kamionkowski,
``Cosmological signature of new parity-violating interactions,''
Phys.\ Rev.\ Lett.\ {\bf 83} (1999) 1506
[astro-ph/9812088].

\bibitem{Finelli2009}
F.~Finelli and M.~Galaverni,
``Rotation of linear polarization plane and circular polarization from cosmological pseudoscalar fields,''
Phys.\ Rev.\ D {\bf 79} (2009) 063002
[arXiv:0802.4210].

\bibitem{MinamiKomatsu2020}
Y.~Minami and E.~Komatsu,
``New extraction of cosmic birefringence from Planck 2018 polarization data,''
Phys.\ Rev.\ Lett.\ {\bf 125} (2020) 221301
[arXiv:2011.11254].

\bibitem{Kamionkowski2009}
M.~Kamionkowski,
``How to De-Rotate the Cosmic Microwave Background Polarization,''
Phys.\ Rev.\ Lett.\ {\bf 102} (2009) 111302
[arXiv:0810.1286].

\bibitem{Planck2018}
Planck Collaboration,
``Planck 2018 results. XI. Polarized dust foregrounds,''
Astron.\ Astrophys.\ {\bf 641} (2020) A11
[arXiv:1801.04945].

\bibitem{ACT2023}
ACT Collaboration,
``Detection of cosmic birefringence from CMB polarization data,''
[JCAP or Phys.\ Rev.\ reference, with arXiv].

\bibitem{Weinberg2008}
S.~Weinberg,
\textit{Cosmology},
Oxford University Press (2008).

\bibitem{Mukhanov2005}
V.~Mukhanov,
\textit{Physical Foundations of Cosmology},
Cambridge University Press (2005).

\bibitem{Dodelson2003}
S.~Dodelson,
\textit{Modern Cosmology},
Academic Press (2003).

\bibitem{Ising1925}
E.~Ising,
``Beitrag zur Theorie des Ferromagnetismus,''
Z.\ Phys.\ {\bf 31} (1925) 253.

\bibitem{Onsager1944}
L.~Onsager,
``Crystal statistics. I. A two-dimensional model with an order-disorder transition,''
Phys.\ Rev.\ {\bf 65} (1944) 117.

\bibitem{KosterlitzThouless1973}
J.~M.~Kosterlitz and D.~J.~Thouless,
``Ordering, metastability and phase transitions in two-dimensional systems,''
J.\ Phys.\ C {\bf 6} (1973) 1181.

\bibitem{Kuramoto1984}
Y.~Kuramoto,
\textit{Chemical Oscillations, Waves, and Turbulence},
Springer (1984).

\bibitem{Strogatz2000}
S.~H.~Strogatz,
``From Kuramoto to Crawford: exploring the onset of synchronization in populations of coupled oscillators,''
Physica D {\bf 143} (2000) 1.

\bibitem{Goldenfeld1992}
N.~Goldenfeld,
\textit{Lectures on Phase Transitions and the Renormalization Group},
Addison-Wesley (1992).

\bibitem{Kubo1966}
R.~Kubo,
``The fluctuation-dissipation theorem,''
Rep.\ Prog.\ Phys.\ {\bf 29} (1966) 255.

\bibitem{Seifert2012}
U.~Seifert,
``Stochastic thermodynamics, fluctuation theorems and molecular machines,''
Rep.\ Prog.\ Phys.\ {\bf 75} (2012) 126001.

\bibitem{Ozawa2019}
M.~Ozawa et al.,
``Topological photonics,''
Rev.\ Mod.\ Phys.\ {\bf 91} (2019) 015006.

\bibitem{Joannopoulos2008}
J.~D.~Joannopoulos, S.~G.~Johnson, J.~N.~Winn and R.~D.~Meade,
\textit{Photonic Crystals: Molding the Flow of Light},
Princeton University Press (2008).

\end{thebibliography}

\end{document}

